%% This is a JAIR Example File Compiled by Nicholas Mattei (nsmattei@tulane.edu) 
%% and Odd Erik Gundersen (odderik@ntnu.no)
%% and Mykel Kochenderfer (mykel@stanford.edu)
%% March 29, 2025
%%
%% This file is based off the ACM Latex Template https://www.acm.org/publications/proceedings-template
%% Revision 2.12 (12/28/2024)
%% 
%% Please see https://www.jair.org/index.php/jair for more information and submission instructions.
%%

%% The first command in your LaTeX source must be the \documentclass
%% command.
%%
%% For submission and review of your manuscript please change the
%% command to \documentclass[manuscript, screen, review]{jair}.
%%
%\documentclass[manuscript, screen, review]{jair}
\documentclass[review]{jair}

\usepackage{latexsym}
\let\Bbbk\relax
\usepackage{amssymb}
\usepackage{amsmath}
\usepackage{amsthm}
\usepackage{booktabs}
\usepackage{enumitem}
\usepackage{graphicx}
\usepackage{color}
%\usepackage{algorithm}
\usepackage{algpseudocode}
\usepackage{multirow}
\usepackage{subcaption}
\usepackage{caption}
\newtheorem{theorem}{Theorem}
\newtheorem{lemma}[theorem]{Lemma}
\newtheorem{corollary}[theorem]{Corollary}
\newtheorem{proposition}[theorem]{Proposition}
\newtheorem{fact}[theorem]{Fact}
\newtheorem{definition}{Definition}
\usepackage[linesnumbered,ruled,vlined]{algorithm2e}


\DeclareMathOperator*{\argmax}{arg\,max}
\DeclareMathOperator*{\argmin}{arg\,min}
\setcopyright{cc}
\copyrightyear{2025}
\acmYear{2025}
\acmDOI{10.1613/jair.1.xxxxx}

%%
\JAIRAE{Daniel Harabor and Sven Koenig}
\JAIRTrack{Multi-Agent Path Finding}
\acmVolume{4}
\acmArticle{111}
\acmMonth{5}
\acmYear{2025}

%%
%% For managing citations, it is recommended to use bibliography
%% files in BibTeX format.
%%
%% You can then either use BibTeX with the ACM-Reference-Format style,
%% or BibLaTeX with the acmnumeric or acmauthoryear sytles, that include
%% support for advanced citation of software artefact from the
%% biblatex-software package, also separately available on CTAN.
%%
%% Look at the sample-*-biblatex.tex files for templates showcasing
%% the biblatex styles.
%%

%%
%% The majority of ACM publications use numbered citations and
%% references.  The command \citestyle{authoryear} switches to the
%% "author year" style.
%%
%% If you are preparing content for an event
%% sponsored by ACM SIGGRAPH, you must use the "author year" style of
%% citations and references.
%% Uncommenting
%% the next command will enable that style.
%%\citestyle{acmauthoryear}


%%
%% end of the preamble, start of the body of the document source.
\begin{document}

%%
%% The "title" command has an optional parameter,
%% allowing the author to define a "short title" to be used in page headers.
%\title{JAIR Example Template}
\title[Adaptive Neighborhood Selection for MAPF via Non-Stationary Bandits]{Adaptive Neighborhood Selection for MAPF via Non-Stationary Bandits}

%%
%% The "author" command and its associated commands are used to define
%% the authors and their affiliations.
%% Of note is the shared affiliation of the first two authors, and the
%% "authornote" and "authornotemark" commands
%% used to denote shared contribution to the research and/or corresponding author.
\author{Amath Sow}
\authornote{Corresponding Author.}
\email{amath.sow@liu.se}
\orcid{0009-0006-5681-9797}
%\orcid{https://orcid.org/0000-0002-2037-3694} % This works too if you want to see the full URL
\affiliation{%
  \institution{Linköping University}
  \city{Linköping}
  \state{Östergötland}
  \country{Sweden}
}

\author{Daniel de Leng}
\orcid{0000-0001-6356-045X}
\email{daniel.de.leng@liu.se}
\affiliation{%
 \institution{Linköping University}
  \city{Linköping}
  \state{Östergötland}
  \country{Sweden}}

\author{Mariusz Wzorek}
\orcid{0000-0003-2147-2114}
\email{mariusz.wzorek@liu.se}
\affiliation{%
  \institution{Linköping University}
  \city{Linköping}
  \state{Östergötland}
  \country{Sweden}
}

\author{Fredrik Heintz}
\orcid{0000-0002-9595-2471}
\email{fredrik.heintz@liu.se}
\affiliation{%
  \institution{Linköping University}
  \city{Linköping}
  \state{Östergötland}
  \country{Sweden}
}


%%
%% By default, the full list of authors will be used in the page
%% headers. Often, this list is too long, and will overlap
%% other information printed in the page headers. This command allows
%% the author to define a more concise list
%% of authors' names for this purpose.
%\renewcommand{\shortauthors}{Trovato et al.}

%%
%% The abstract is a short summary of the work to be presented in the
%% article.
\begin{abstract}
{\bf Abstract:}
Adaptive Large Neighborhood Search (ALNS) is a meta-heuristic framework widely used to solve large-scale MAPF problems. It uses selection mechanisms such as the Roulette Wheel or Multi-Armed Bandit (MAB) policies, which adaptively learn to identify the most effective neighborhood (heuristic) at each iteration. A key assumption of these policies is that each heuristic's reward distribution (effectiveness) is stationary. However, the best heuristic can change during the search process. In this work, we address the inherent non-stationarity of neighborhood selection in ALNS-based MAPF solvers, where the reward distribution of the heuristic changes over time. We introduce DyMAB (Dynamic multi-armed bandit for neighborhood selection), a novel anytime algorithm that integrates a non-stationary multi-armed bandit (MAB) framework into ALNS.
DyMAB leverages a windowed queue to track recent rewards, allowing for dynamic adaptation of neighborhood selection at each iteration. This ensures that the most effective heuristic is prioritized, improving overall search efficiency. We present a theoretical analysis showing that DyMAB achieves sub-linear regret under standard non-stationary bandit assumptions, offering formal performance guarantees.
We benchmark DyMAB against a range of state-of-the-art anytime solvers across diverse MAPF scenarios, demonstrating that our method significantly enhances solution quality and scalability. In most cases, DyMAB achieves up to a $50\%$ reduction in AUC and up to a $30\%$ decrease in the Sum of Delays, highlighting its superiority in optimizing large-scale MAPF instances.
\end{abstract}


%% JAIR Note: 
%% Do not include ACM CCS Concepts or Keywords


%% To be updated by authors.
%\received{20 February 2007}
%\received[revised]{12 March 2009}
%\received[accepted]{5 June 2009}

%%
%% This command processes the author and affiliation and title
%% information and builds the first part of the formatted document.
\maketitle

\section{Introduction}
\text{Multi-Agent Path Finding (MAPF)}~\citep{SternEtal2019} is a problem where multiple agents navigate in a shared environment from their start locations to goal locations without collisions. Different solvers have been proposed with varying levels of success, broadly classified as \emph{optimal} and \emph{near-optimal} approaches~\citep{SBOS}. Optimal solvers, such as CBS~\citep{CBS}, guarantee that a least-cost solution will be found using a two-level search strategy. Still, their scalability is limited to small instances of problems due to exponential time complexity. To address scalability, near-optimal solvers trade-off solution quality for faster computations, offering practical alternatives for large and complex environments~\citep{ECBS, EECBS}.
Among near-optimal solvers, anytime algorithms have become very popular for their focus on finding near-optimal solutions within a given time budget. MAPF-LNS~\citep{LNS}, a state-of-the-art anytime solver, uses the Large Neighborhood Search (LNS) meta-heuristic technique from combinatorial optimization~\citep{ori-LNS, LNS-comb}. This method iteratively refines an initial sub-optimal solution by employing destroy heuristics to remove paths of selected agents and repair operators to replan these paths, aiming to minimize the overall cost. ALNS extends LNS by maintaining a set of heuristics. Instead of using a fixed strategy, ALNS aims to adaptively select the most promising heuristic at each iteration based on their observed performance. Techniques like Roulette Wheel~\citep{LNS} or Multi-Armed Bandit (MAB)~\citep{Balance} have been employed for this adaptive selection mechanism. 

However, a fundamental assumption underlying these adaptive selection mechanisms is the stationarity of the reward distribution associated with each heuristic. This assumption means that the effectiveness of a particular heuristic remains constant during the search. In the context of MAPF, this may not accurately reflect the dynamic nature of the solution space. More specifically, different heuristics target different aspects of the solution space. For example, some focus on the most delayed agents, others on collision graphs, and some target dense environmental bottlenecks. As the solution improves and conflicts are resolved, the performance of these heuristics can shift significantly. Consequently, a heuristic that was effective early in the search may become less performant later on. Therefore, relying on historical performance under a stationarity assumption may lead to suboptimal heuristic selection and hinder the overall search efficiency.

\begin{figure*}[t]
    \centering
    \includegraphics[width=\linewidth]{./images/DyMAB2.png}
    \caption{Overview of DyMAB algorithm. An initial sub-optimal solution is generated using Prioritized Planning. $H_1$, $H_2$ and $H_3$ constitute set of heuristics for DyMAB. At each step, we have $r_{i\_{H_i}}$, which is the reward after applying the heuristic $H_i$. We maintain a reward memory $Q$ of size $W$ to store only the most recent rewards for each heuristic $H_i$. A heuristic selection step occurs. This step chooses which heuristic ($H_1, H_2, H_3$) to apply next. The policy aims to balance exploring different heuristics and exploiting the one that has recently performed best based on the policy used, such as DyMAB($\epsilon$-Greedy) and DyMAB($\alpha$-UCB) with logarithmic decay. The solution exploration continues, i.e., the process of modifying the current solution (destroying a subset of the solution) based on the chosen heuristic and replanning the paths in order to improve the overall solution cost. The process continues until the time budget is exhausted. }
    \label{fig:dymab}
\end{figure*}

To address this limitation, we introduce \text{DyMAB} (Dynamic Multi-Armed Bandit for Neighborhood Selection), a novel anytime MAPF solver specifically designed to address the inherent non-stationarity of heuristic performance within the ALNS framework. The DyMAB architecture is illustrated in Fig~\ref{fig:dymab}.
The key contributions of our work are 1) two novel destroy heuristics specifically designed for ALNS-based solvers; 2) two novel MAB policies, \text{DyMAB}($\alpha$-UCB) and \text{DyMAB}($\epsilon$-Greedy), explicitly designed for non-stationary MAB; and 3) a theoretical analysis of the DyMAB framework, establishing dynamic regret bounds and optimal parameter conditions under piecewise-stationary reward assumptions.

The remainder of this paper is organized as follows.
An overview of related work is presented in Section~\ref{sec:related-work}.
Section~\ref{sec:background} introduces the MAPF problem and its relation to non-stationary multi-armed bandits.
Section~\ref{sec:dymab} introduces our \text{DyMAB} approach.
Section~\ref{sec:theory} presents theoretical analysis of \text{DyMAB}.
An experimental evaluation showing the performance of \text{DyMAB} relative to the state-of-the-art is presented in Section~\ref{sec:experiments}.
The paper ends with conclusions and a discussion of future work in Section~\ref{sec:conclusion}.

\begin{comment}
    


\section{Acknowledgments}

Identification of funding sources and other support, and thanks to
individuals and groups that assisted in the research and the
preparation of the work should be included in an acknowledgment
section, which is placed just before the reference section in your
document.

This section has a special environment:
\begin{verbatim}
  \begin{acks}
  ...
  \end{acks}
\end{verbatim}
so that the information contained therein can be more easily collected
during the article metadata extraction phase, and to ensure
consistency in the spelling of the section heading.

Authors should not prepare this section as a numbered or unnumbered {\verb|\section|}; please use the ``{\verb|acks|}'' environment.
\end{comment}


%%%%%%%%%%%%%%%%%%%%%%%%%%%%%%%%%%%%%%%%%%%%%%%%%%%%%%%%%%%%%%%%%%%%%%%%
%%%%%%%%%%%%%%%%%%%%%%%%%%%%%%%%%%%%%%%%%%%%%%%%%%%%%%%%%%%%%%%%%%%%%%%%
\section{Related work} \label{sec:related-work}
Anytime MAPF solvers are a subset of sub-optimal solvers that focus on finding near-optimal solutions within a time budget, enabling a trade-off between time and solution quality. One of the leading approaches in this category is MAPF-LNS~\citep{LNS}, which has been applied effectively to large-scale MAPF problems. It uses a sub-optimal solver to generate an initial solution~\citep{PP, PPS} and then iteratively improves the solution. This is done by removing paths for a subset of agents (destroy heuristic), replanning these paths (repair operator), and selecting the solution with the lowest cost. The process continues until the time budget is exhausted. Adaptive LNS (ALNS) further enhances this approach by using a set of competing heuristics based on their historical performance. These destroy heuristics are a random uniform selection of N paths, an agent-based heuristic ($H_{\text{agent}}$), and a map-based heuristic ($H_{\text{grid}}$)~\citep{LNS}. The agent-based heuristic $H_{\text{agent}}$ focuses on the path of the agent with the maximum delay and other paths that affect it. The map-based heuristic $H_{\text{grid}}$ randomly chooses a vertex with a high degree in the graph, which includes paths through the selected vertex. Then, a selection policy $\pi$ such as a Roulette Wheel is used to choose between the set of heuristics $H$ by maintaining updatable weights $w$ \citep{ori-LNS}.
Recent advancements include Machine Learning-Guided LNS (ML-LNS)~\citep{ML-LNS}, which uses imitation learning to guide the selection of agent paths for replanning, and \text{MAPF-LNS2}~\citep{LNS2}, which integrates SIPPS algorithm~\citep{SIPP} for faster single-agent path-finding, significantly speeding up the solution process and allowing for efficient solving of large scenarios with thousands of agents.

A notable recent algorithm is \text{BALANCE} (Bandit-based Adaptive Large Neighborhood Search Combined with Exploration)~\citep{Balance}. BALANCE uses a bi-level MAB scheme to adapt the selection of destroy heuristics and neighborhood sizes. It uses policies based on \text{Upper Confidence Bound (UCB1)}, and \text{Thompson Sampling} that outperform previous \text{MAPF-LNS} solvers. However, these policies assume that the performance of the heuristics is stationary. In this paper, we demonstrate that removing this assumption enables more effective neighborhood selection, significantly improving solution costs and search efficiency even in large-scale scenarios with more than thousand agents. ADDRESS~\citep{ADDRESS} is another recent anytime LNS-based solver that employs Thompson Sampling with a diverse pool of destroy heuristics, achieving strong anytime convergence performance. LaCAM~\citep{LaCam} is a scalable complete MAPF solver based on lazy constraints, and LaCAM*~\citep{LaCamStar} extends it into an eventually optimal solver by iteratively refining solutions after an initial one is found. 



%%%%%%%%%%%%%%%%%%%%%%%%%%%%%%%%%%%%%%%%%%%%%%%%%%%%%%%%%%%%%%%%%%%%%%%%%%%%%%%%%%%%%%%%%%%%%%%%%%%%%%%%%%%%%%%%%%%%
\section{Background} \label{sec:background}
\subsection{Problem formalization}
We consider a MAPF problem~\citep{SternEtal2019}, which is modeled as an undirected graph $G = (V, E)$, where $V$ is the set of vertices and $E$ is the set of edges. The problem involves $N$ agents, denoted as $A = \{ a_1, a_2, \dots, a_N \}$. 
Each agent $a_i \in A$ starts at vertex $s_i \in V$ and aims to reach a goal vertex $g_i \in V$ by executing a sequence of actions over discrete time-steps. At each time-step, an agent can move to an adjacent vertex or wait at its current vertex. A path \( \mathcal{P}_i \) for agent \( a_i \) is a sequence of vertices: \(\mathcal{P}_i = [p_{i_0}, p_{i_1}, \dots, p_{i_{l(\mathcal{P}_i)}}]\), where  \( p_{i_0} = s_i \) ,  \( p_{i_{l(\mathcal{P}_i)}} = g_i \), and \( l(\mathcal{P}_i) \) is the length or travel distance of the path \citep{LNS}.
We consider two types of conflicts. A Vertex conflict $\left< a_i, a_j, v, t \right>$ occurs when two agents $a_i$ and $a_j$ attempt to occupy the same vertex $v$ at the same time $t$ and an Edge conflict $\left< a_i, a_j, u, v, t \right>$ occurs when the agents attempt to traverse the same edge $(u,v) \in E$ in opposite directions during the same time $t$ \citep{Stern2019}.

A valid solution is a set of collision-free paths $\mathcal{P} = \{\mathcal{P}_i \mid a_i \in A \}$. The objective is to find a solution $\mathcal{P}$ that minimizes the \text{Sum of Costs (SOC)}, i.e., the total travel distance of all agents, defined as:  
\begin{equation}
    \argmin_\mathcal{P} \sum_{\mathcal{P}_i \in \mathcal{P}} l(\mathcal{P}_i).
\end{equation}

In the Anytime MAPF framework, the goal is equivalently expressed as minimizing the Sum of Delays~\citep{LNS}:
\begin{equation}
\argmin_\mathcal{P} \sum_{\mathcal{P}_i \in \mathcal{P}} delays(\mathcal{P}_i),
\end{equation}
where the delay of an agent $a_i$'s path is computed as:
\begin{equation}
delays(\mathcal{P}_i) = l(\mathcal{P}_i) - d(s_i, g_i),
\end{equation}
with $d(s_i, g_i)$ being the shortest path distance (in terms of graph edges) between the start vertex $s_i$ and the goal vertex $g_i$.

Another important metric for evaluating the performance of anytime solvers is the Area Under the Curve (AUC), which represents the cumulative Sum of Delays (SOD) 
\begin{equation}
\sum_{\mathcal{P}_i \in \mathcal{P}} delays(\mathcal{P}_i),
\end{equation}
as a function of runtime. The AUC is defined as:
\begin{equation}
AUC(t_{\text{limit}}) = \int_{t_{\text{init}}}^{t_{\text{limit}}} SOD(t) \, dt,
\end{equation}
where a lower AUC indicates faster improvement in solution quality~\citep{ML-LNS}.

Another standard metric in MAPF is the \text{Makespan}, defined as $\max_i\, l(\mathcal{P}_i)$, i.e., the time at which the last agent reaches its goal. In this work, we focus on SOC/SOD minimization, which is the standard objective for LNS-based solvers~\citep{LNS, LNS2}.


\subsection{Non-stationary Multi-Armed Bandits}
A $K$-armed bandit is a sequential decision problem in which, at each round $t \in \{1, \ldots, T\}$, an agent selects an arm $k_t \in \{1, \ldots, K\}$ and receives a reward $r_{k_t}(t)$ drawn from a distribution with unknown mean $\mu_{k_t}$. The goal is to maximize the cumulative reward $\sum_{t=1}^T r_{k_t}(t)$, or equivalently to minimize the cumulative regret $R_T = \sum_{t=1}^T (\mu^* - \mu_{k_t})$, where $\mu^* = \max_k \mu_k$.

In ALNS-based MAPF solvers, the algorithm's performance critically depends on selecting appropriate destroy heuristics at each iteration. Given multiple competing heuristics (each with unknown and evolving performance), how can we decide which one to apply at each step to maximize overall solution improvement? This naturally aligns with the MAB framework, where each heuristic corresponds to an arm and pulling an arm yields a reward representing the improvement in solution cost. While classical MAB methods have been applied to this setting, they assume stationary reward distributions, which fail to capture the evolving nature of heuristic performance during the search.

The non-stationary MAB problem extends the classic MAB formulation by allowing reward distributions to change over time. At each time step \(t\), the agent selects an arm \(k\) from a set of \(K\) arms and receives a reward \(r_k(t)\), drawn from a time-varying distribution with a mean of \(\mu_k(t)\)~\citep{MAB1, Non-Stationary1}. The objective is to maximize the cumulative reward, \(\sum_{t=1}^{T} r_k(t)\), over a time horizon \(T\). This setup introduces a unique challenge: balancing exploration (trying new or less certain arms) and exploitation (choosing arms known to yield high rewards) while adapting to temporal variability in reward distributions.
Several approaches have been proposed to address non-stationarity. \text{Sliding-Window UCB (SW-UCB)}~\citep{Non-Stationary} adapts traditional UCB methods by considering only the rewards obtained within a fixed sliding window. This allows the algorithm to prioritize recent rewards, making it responsive to changes in reward distributions. In contrast, \text{Discounted UCB (D-UCB)}~\citep{Non-Stationary2} uses a discount factor \(\gamma \in (0,1]\) to reduce the weight of older rewards exponentially. This method retains a memory of past observations but gradually diminishes their influence, increasing sensitivity to recent changes. Similar strategies have been applied to Thompson Sampling~\citep{windowed-thompson, discounted-thompson}.  

%%%%%%%%%%%%%%%%%%%%%%%%%%%%%%%%%%%%%%%%%%%%%%%%%%%%%%%%%%%%%%%%%%%%%%%%
\section{DyMAB for MAPF} \label{sec:dymab}
To enable adaptive heuristic selection within the ALNS framework for MAPF, we propose DyMAB, a MAB approach that accounts for the non-stationary performance of heuristics during search. Specifically, we introduce two non-stationary MAB policies tailored for this setting: \text{DyMAB}($\alpha$-UCB) and \text{DyMAB}($\epsilon$-Greedy). These policies maintain a sliding reward window queue $Q$ of size \(W\) for each heuristic to prioritize recent performance and apply a \emph{logarithmic decay} to the exploration parameters \(\alpha\) and \(\epsilon\), encouraging early exploration followed by gradual exploitation of the most effective heuristics. Alternative formulations could be considered — for instance, modeling the full interaction as an MDP, or using a contextual bandit that conditions selection on features of the current solution state. We leave that as future work.In this section, we detail the neighborhood selection strategies and the DyMAB policies used to guide their application.


\subsection{Neighborhood Selection}
Efficient exploration of the solution space in ALNS relies on well-defined heuristics. This section introduces three distinct heuristics, each tailored to explore the solution space. 

\subsubsection{Random Neighborhood}
For the \emph{random neighborhood heuristic} $H_{\textnormal{rand}}$, we randomly select uniformly a subset of agents \(A_s\) given a neighborhood size \(M\). It is a naive approach and has shown interesting results, particularly for congested environments~\citep{Emrah, LNS, Song, LNS2}.

\subsubsection{Synchronized Random Walk Neighborhood} 
The \emph{Synchronized random walk neighborhood} $H_{\textnormal{sync}}$ is inspired by the agent-based neighborhood approach presented in the \text{MAPF-LNS} framework~\citep{LNS}. However, unlike the original approach, which selects only the most delayed agent, $H_{\textnormal{sync}}$ generalizes this by defining a set of seed agents \( S \) based on a delay metric \( D \) and then simultaneously performing random walk. This method balances between the most delayed agents and those blocking their progress, which can lead to more effective path repair and better overall solutions. 

\begin{comment}  
Specifically, \( S = \{a_1, a_2, \ldots, a_k\} \), where  \( k \) is randomly sampled from the range \([2, round(M/2)]\) and \( M > 2 \) represents the desired neighbourhood size. The agents \( a_i \) in \( S \) are the \( k \)  most delayed agents in the initial solution generated. This method balances between the most delayed agents and those blocking their progress, which can lead to more effective path repair and better overall solutions. 

The initial neighbourhood \( \mathcal{N} \) is initialized to include all seed agents: \( \mathcal{N} = \{a_1, a_2, \ldots, a_k\} \). Starting with these seeds, \( \mathcal{N} \) is iteratively expanded using a synchronised random walk, aiming to identify agents whose paths conflict with the seeds.
For each seed agent \( a_i \), its current state at time \( t \) is defined as \( \text{current\_state}(a_i, t) = (x_i(t), t) \), where $x_i(t)$ is the position vertex of $a_i$ at time $t$. From its next valid locations, including the option to remain in place, we randomly select a move as follows:
\begin{equation}
x_i(t+1) \sim \text{Uniform}(\text{next\_locations}(a_i, t)),
\end{equation}
where
\begin{equation}
\text{next\_locations}(a_i, t) = \left\{ x \,\middle|\,
\begin{aligned}
& (x_i(t), x) \in E \cup \{x_i(t)\} \ \text{and} \\
& t+1+h(x) < l(p_i)
\end{aligned}
\right\}.
\end{equation}
Here, \( h(x) \) denotes the heuristic estimate of the cost from \( x \) to the goal, and \( l(p_i) \) is the length of agent \( a_i\)'s current path. This feasibility constraint encourages movements that remain within the agent's planning horizon and search toward goal-directed progress.

Conflicts are checked after each move, and any newly identified conflicting agents are added to \( \mathcal{N} \). The synchronised random walk continues until the neighbourhood \( \mathcal{N} \) reaches the desired size \( M \).
\end{comment}
The corresponding pseudo-code is presented in Algorithm~\ref{alg:srwn}.
The algorithm begins by selecting a set of seed agents \( S \) (line~3), sampled from the top-\( k \) most delayed agents that are not currently in the tabu list \( \mathcal{T} \), where \( k \) is randomly drawn from the interval \([2, \lfloor M/2 \rfloor]\) (line~2). These agents form the initial neighborhood \( \mathcal{N} \) (line~4) and are also added to the tabu list to avoid repeated selection. Then, a synchronized random walk is performed. Each seed agent \( a_i \in S \) is initialized at the start of its path, with \( x_i \) representing its current location and \( t_i = 0 \) as the starting timestep (line~7). The upper bound \( ub_i \) is set to the length of its current path minus one, restricting how far in timesteps the agent is allowed to walk (line~8). We also initialize the variables \texttt{steps}, which tracks the number of walk iterations, and \texttt{no\_progress}, which counts the number of steps taken without any change in the neighborhood \( \mathcal{N} \) (line~9).
At the beginning of each iteration, the algorithm stores the current neighborhood size (line~11) to later detect whether progress has been made. Then, for each seed agent \( a_i \) (line~12), it identifies a set of feasible next locations \( N_{x_i} \) (line~13), consisting of all neighboring vertices (including the current location) such that moving there would allow the agent to still reach its goal within the length of its planned path. This feasibility check ensures that the random walk does not unnecessarily extend the agent’s path beyond its original upper bound. This is enforced through the feasibility constraint: \(t_i + 1 + h(x) < l(\mathcal{P}_i),\) where \( h(x) \) is a heuristic estimate of the remaining cost to the goal.
A random feasible location \( y \in N_{x_i} \) is then selected for agent \( a_i \). All agents that would conflict with \( a_i \) at location \( y \) and time \( t_i + 1 \) are identified and added to the neighborhood \( \mathcal{N} \) via set union (line~14), so duplicates across multiple seed agents are automatically discarded. If \( |\mathcal{N}| \) exceeds \( M \) after this addition, \( \mathcal{N} \) is immediately truncated to retain exactly \( M \) agents (line~15). The walk state of \( a_i \) is then updated to the new position \( y \) and timestep \( t_i + 1 \), and the inner loop breaks (lines~16--17).

After processing all seed agents, the algorithm checks whether the neighborhood size has increased compared to the previous iteration. If \( |\mathcal{N}| > \texttt{prev} \), the counter \texttt{no\_progress} is reset to zero, indicating that progress has been made. Otherwise, \texttt{no\_progress} is incremented to track stagnation (lines 19--22). 
If no progress has been made for \( \lfloor \texttt{max} / 4 \rfloor \) consecutive steps, each seed agent is randomly reset to a different location along its current path (lines 23--26). This re-randomization serves to reinitialize the synchronized walk and helps the algorithm escape potential local stagnation. The global step counter \texttt{steps} is incremented (line 27).
Once the main loop terminates—either because the neighborhood \( \mathcal{N} \) has reached the target size \( M \) or the step limit has been exceeded—the constructed neighborhood \( \mathcal{N} \) is returned (line 28).

\begin{algorithm}[t]
\caption{Synchronized Random Walk Neighborhood}
\label{alg:srwn}
\KwIn{Graph $G = (V, E)$, agents $A = \{a_1, \dots, a_m\}$, paths $P = \{p_1, \dots, p_m\}$, neighborhood size $M$, tabu list $\mathcal{T}$}
\KwOut{Neighborhood $\mathcal{N} \subseteq A$}

\tcc{Step 1: Seed agent selection}
$k \gets$ \textsc{RandInt}$(2, \lfloor M/2 \rfloor)$\;
$S \gets$ Top-$k$ delayed agents from $A \setminus \mathcal{T}$\;
$\mathcal{N} \gets S$; \quad $\mathcal{T} \gets \mathcal{T} \cup S$\;

\tcc{Step 2: Synchronized random walk}
\ForEach{$a_i \in S$}{
  $(x_i, t_i) \gets (p_{i_0}.\text{loc}, 0)$\;
  $ub_i \gets l(\mathcal{P}_i) - 1$\;
}
$steps \gets 0$; \quad $max \gets |A|$; \quad $no\_progress \gets 0$\;

\While{$steps < max \land |\mathcal{N}| < M$}{
  $prev \gets |\mathcal{N}|$\;

  \ForEach{$a_i \in S$}{
    $N_{x_i} \gets \{x \in \text{neighbors}(x_i) \cup \{x_i\} \mid t_i + 1 + h(x) < ub_i\}$\;

    \While{$N_{x_i} \neq \emptyset$ $\land$ $|\mathcal{N}| < M$}{
      $y \gets$ random $x \in N_{x_i}$\;
      $\mathcal{N} \gets \mathcal{N} \cup$ \textsc{ConflictingAgents}$(a_i, x_i, y, t_i + 1)$\;
      \If{$|\mathcal{N}| > M$}{$\mathcal{N} \gets \mathcal{N}[{:}M]$\;}
      $(x_i, t_i) \gets (y, t_i + 1)$\;
      \textbf{break}\;
    }
  }

  \eIf{$|\mathcal{N}| > prev$}{
    $no\_progress \gets 0$\;
  }{
    $no\_progress \gets no\_progress + 1$\;
    \If{$no\_progress \geq \lfloor max / 4 \rfloor$}{
      \ForEach{$a_i \in S$ \textbf{with} $ub_i > 1$}{
        $t_i \gets$ \textsc{RandInt}$(0, ub_i)$\;
        $x_i \gets p_{i_{t_i}}.\text{loc}$\;
      }
    }
  }
  $steps \gets steps + 1$\;
}
\Return $\mathcal{N}$\;
\end{algorithm}




\subsubsection{Entropy-based Neighborhood}
The \emph{Entropy-based neighborhood heuristic} $H_{\textnormal{entr}}$ leverages Shannon information theory, commonly applied in Vehicle Routing Problems (VRP) \citep{entropy2, entropy1}, to identify agents with high path entropy. In our case, path entropy quantifies the variability of an agent's movements, considering both the frequency of location visits in its precomputed path and the likelihood of an agent visiting a particular location at each timestep. The key intuition is to prioritize agents with highly complex paths, as these are often found in congested areas where conflicts are more likely. Resolving such conflicts can significantly enhance the overall solution in MAPF.

\begin{definition}[Path entropy]
Let \( \mathcal{P}_i = [l_0, l_1, \ldots, l_k] \) represent the sequence of locations (vertices) in agent \( a_i \)'s path, and \( \text{freq}(l) \) denote the frequency of location \( l \in V\) in the path, $l_0=p_{i_0}$, $l_1=p_{i_1}$ and $l_k=p_{i_{l(\mathcal{P}_i)}}$. The path entropy of agent \( a_i \), denoted as \( H(\mathcal{P}_i) \), is calculated using Shannon entropy:

\begin{equation}
H(\mathcal{P}_i) = -\sum_{l \in \mathcal{P}_i} p(l) \log_2 p(l),
\end{equation}
where \( p(l) = \text{freq}(l)/|\mathcal{P}_i| \) is the probability of visiting location \( l \).
We define an entropy score \( E(a_i) \) for each agent as:
\begin{equation}
E(a_i) = H(\mathcal{P}_i) \cdot (1 + D(a_i)),
\end{equation}
where \( D(a_i) = l(\mathcal{P}_i) - d(s_i, g_i) \) is the delay factor for agent \( a_i \), consistent with $\mathit{delays}(\mathcal{P}_i)$ as defined in Section~\ref{sec:background}.
\end{definition}

To select the seed agent, we use Boltzmann Sampling, a probabilistic technique that ensures diverse selection while favoring agents with higher entropy scores. The probability of selecting agent \( a_i \) is computed as:
\begin{equation}
p_b(a_i) = \frac{\exp\left(\frac{E(a_i)}{T}\right)}{\sum_{a \in A} \exp\left(\frac{E(a)}{T}\right)},
\end{equation}
where \( T \) is the temperature parameter controlling the randomness of selection. Higher \( T \) values promote exploration by increasing the likelihood of selecting agents with lower scores.
%
The seed agent \( a_i \) is chosen such that:
\begin{equation}
r \leq \sum_{j=1}^i p_b(a_j),
\end{equation}
where \( r \sim U[0, 1] \) is a uniform random value.

After selecting the seed agent, an entropy-guided walk is performed to expand the neighborhood \( \mathcal{N} \). At each timestep \( t \), for every potential next location \( l \), a score \( E(l) \) is computed as:

\begin{equation}
E(l) = \frac{H(l, t)}{1 + \text{h}(l)},
\end{equation}
where
\( H(l, t) \) is the location entropy at \( (l, t) \).

\begin{definition}[Location entropy]
The location entropy is calculated as
\begin{equation}
H(l, t) = -\sum_{a_i \in A} p(a_i) \log_2 p(a_i),
\end{equation}
where
\begin{equation}
p(a_i) = \frac{\text{freq}(a_i, l, t)}{\sum_{a_j \in A} \text{freq}(a_j, l, t)}
\end{equation}
represents the likelihood of agent $a_i$ visiting location $l$ at timestep $t$, and $\text{freq}(a_i, l, t) = \mathbf{1}[p_{i_t} = l]$ is the indicator function over \emph{all} agents $a_i \in A$ using their current solution paths. Thus $H(l,t)$ captures the congestion at location $l$ at time $t$ across the entire current solution; cells visited by no agent have zero entropy.
\end{definition}
%

The corresponding pseudo-code is presented in Algorithm~\ref{alg:entropy-neighborhood}. The process begins by computing the path entropy for each agent \( a_i \in A \setminus \mathcal{T} \), where \( \mathcal{T} \) is the tabu list and their entropy score \( E(a_i) \) (lines~2--4). 
Then, a Boltzmann distribution is computed over the entropy scores and normalized. Each agent's probability of being selected is proportional to the exponential of its entropy score, scaled by a temperature parameter \( T \) that controls the degree of exploration(lines~5--13). 
A single seed agent is then sampled according to this distribution. This is done by generating a uniform random number \( r \in [0, 1] \) and selecting the first agent \( a_i \) such that the cumulative probability up to \( a_i \) exceeds \( r \) (lines~14--21).
The selected seed agent is then added to the neighborhood \(\mathcal{N}\) (line~23), and an entropy-guided walk is performed to expand the neighborhood (line~24), after which the resulting neighborhood \(\mathcal{N}\) is returned (line~25). 

The entropy-guided walk is presented in Algorithm~\ref{alg:entropy-random-walk}. It starts by initializing the agent's position to the start of its current path and setting the timestep \( t = 0 \) (lines~1--2). At each iteration, the agent evaluates all possible next locations, including wait action (line~3--4). 
For each candidate location \( l \), the algorithm calculates a score based on the location entropy \( H(l, t+1) \), and the heuristic cost to the goal $h(l)$ (lines~5--9).
All candidate moves are sorted in descending order of their scores (line~12), and the agent selects the highest-scoring move that satisfies a temporal feasibility constraint, namely \(t + 1 + h(l) < |\mathcal{P}_i|\) (line~13). Then, any agents that conflict with the move are added to the neighborhood \(\mathcal{N}\) (line~14).

If the agent cannot find any feasible move, the random walk terminates (lines~18--19). Otherwise, the process continues, advancing along the agent’s path, until either the neighborhood reaches the target size \(M\) or the agent finishes traversing its path. The final constructed neighborhood \(\mathcal{N}\) is then returned (line~20).


\begin{algorithm}[t]
\caption{Entropy-Based Neighborhood}
\label{alg:entropy-neighborhood}
\KwIn{Agents $A = \{a_1, \ldots, a_m\}$, Paths $P = \{p_1, \ldots, p_m\}$, Neighborhood size $M$, Tabu list $\mathcal{T}$, Temperature $T$}
\KwOut{Neighborhood $\mathcal{N}$}

\tcc{Step 1: Compute agent path entropies}
\ForEach{$a_i \in A \setminus \mathcal{T}$}{
    $H(\mathcal{P}_i) \gets -\sum_{l \in \mathcal{P}_i} \frac{\text{freq}(l)}{|\mathcal{P}_i|} \log_2 \left( \frac{\text{freq}(l)}{|\mathcal{P}_i|} \right)$\;
    $E(a_i) \gets H(\mathcal{P}_i) \cdot (1 + D(a_i))$\;
}

\tcc{Step 2: Boltzmann probabilities}
$probabilities \gets [\,]$\;
$Z \gets 0$\;

\ForEach{$E(a_i)$}{
    $s_i \gets \exp(E(a_i)/T)$\;
    Append $s_i$ to $probabilities$\;
    $Z \gets Z + s_i$\;
}

\For{$i \gets 1$ \KwTo $|\text{probabilities}|$}{
    $probabilities[i] \gets probabilities[i] / Z$\;
}

\tcc{Step 3: Select seed agent}
$r \gets \textsf{Uniform}(0, 1)$\;
$cumulative\_prob \gets 0$\;

\For{$i \gets 1$ \KwTo $|probabilities|$}{
    $cumulative\_prob \gets cumulative\_prob + probabilities[i]$\;
    \If{$r \leq cumulative\_prob$}{
        $\text{seed\_agent} \gets i$\;
        \textbf{break}\;
    }
}

\tcc{Step 4: Entropy-guided walk}
$\mathcal{N} \gets \{\text{seed\_agent}\}$\;
$\mathcal{N} \gets$ \textsf{LocEntropyRandomWalk}$(\text{seed\_agent}, P, \mathcal{N}, M)$\;

\Return $\mathcal{N}$\;
\end{algorithm}


\begin{algorithm}[t]
\caption{Entropy-Guided Walk}
\label{alg:entropy-random-walk}
\KwIn{Agent $a_i$, Path $\mathcal{P}_i$, Neighborhood $\mathcal{N}$, Neighborhood size $M$}
\KwOut{Updated Neighborhood $\mathcal{N}$}

$t \gets 0$\; 
$loc \gets p_{i_0}.\text{loc}$\;

\While{$t < |\mathcal{P}_i| - 1$ $\land$ $|\mathcal{N}| < M$}{
    $next\_locs \gets \text{neighbors}(loc) \cup \{loc\}$\;
    $location\_scores \gets [\,]$\;

    \ForEach{$l \in next\_locs$}{
        $p(a_j) \gets \frac{\text{freq}(a_j, l, t+1)}{\sum_{a_k \in A} \text{freq}(a_k, l, t+1)}$\;
        $H(l, t+1) \gets -\sum_{a_j \in A} p(a_j) \log_2 p(a_j)$\;
        $E(l) \gets H(l, t+1) \cdot \frac{1}{1 + h(l)}$\;
        Append $(l, E(l))$ to $location\_scores$\;
    }

    $moved \gets \text{false}$\;

    \ForEach{$(l, E(l)) \in \textsf{Sorted}(location\_scores)$}{
        \If{$t + 1 + h(l) < |\mathcal{P}_i|$}{
            $\mathcal{N} \gets \mathcal{N} \cup$ \textsf{ConflictingAgents}$(loc, l, t+1)$\;
            $(loc, t) \gets (l, t+1)$\;
            $moved \gets \text{true}$\;
            \If{$|\mathcal{N}| \geq M$}{\Return $\mathcal{N}$\;}
            \textbf{break}\;
        }
    }

    \If{Not $moved$}{
        \textbf{break}\;
    }
}
\Return $\mathcal{N}$\;
\end{algorithm}


\subsection{DyMAB problem definition}
We formally define \emph{DyMAB} as a non-stationary (i.e., ~dynamic) MAB problem designed to optimize neighborhood selection within ALNS-based MAPF solvers. The core challenge is to effectively manage a set of $K$ arms representing heuristics, $H = \left\{H_{rand}, H_{sync}, H_{entr}\right\}$, where each heuristic $ H_i \in H $ corresponds to a specific neighborhood selection strategy (i.e., random selection, synchronized random walk and entropy-based).

At each time step $ t $, each heuristic $ H_i \in H $ is associated with a reward distribution $ D_i(t) $, characterized by a time-dependent mean reward:
\begin{equation}
\mu_i(t) = \mathbb{E}[r_{i,t}], \quad r_{i,t} \sim D_i(t),
\end{equation}
where $ r_{i,t} $ represents the reward obtained by applying the heuristic $ H_i $ at time $ t $. The time dependency of $ \mu_i(t) $ reflects the non-stationary nature of the reward distributions in DyMAB. 
We first use a sub-optimal algorithm to compute an initial solution $P$. Specifically, we employ \emph{Prioritized Planning} (PP)~\citep{PP}, which is an effective sub-optimal solver as proven in the \text{MAPF-LNS} framework~\citep{LNS}. From the initial solution $ P $, a subset of agents $ A_s \subseteq A $ is selected based on the defined neighborhood. Their paths, denoted as $ P_s^- $, are removed and re-planned to yield new paths $ P_s^+ $. The \emph{reward} for a heuristic \( H_i \) is the improvement in the global SOC objective directly attributable to that heuristic at iteration $t$, expressed as:
\begin{equation}
r_{i,t} = l(\mathcal{P}_s^-) - l(\mathcal{P}_s^+),
\end{equation}
where \( l(\mathcal{P}_s) = \sum_{\mathcal{P}_j \in \mathcal{P}_s} l(\mathcal{P}_j) \) denotes the total path length of the selected agents before ($\mathcal{P}_s^-$) and after ($\mathcal{P}_s^+$) replanning. Since replanning affects only the selected subset and the global SOC objective is $\min \sum_i l(\mathcal{P}_i)$, $r_{i,t}$ is precisely the change in the global objective: a positive value means the solution improved, while zero or negative means it did not.

To address the non-stationary nature of heuristic performance, DyMAB maintains a \emph{windowed reward memory} \( Q_i \) of fixed size \( W \) for each heuristic \( H_i \), inspired by non-stationary bandit frameworks~\citep{Non-Stationary, Non-Stationary1, Non-Stationary2}. After each call of a heuristic, the most recent reward \( r_{i,t} \) is appended to \( Q_i \); when the queue reaches its capacity, the oldest reward is discarded. This sliding window mechanism ensures that the estimate of the expected reward \(\hat{\mu}_i(t)\) for each heuristic is based exclusively on the most recent \( W \) samples:
\begin{equation}
\hat{\mu}_i(t) = \frac{1}{|Q_i|} \sum_{r \in Q_i} r_{i,t}.
\end{equation}
Consequently, this mechanism ensures that the algorithm remains responsive to changes in reward distributions of the heuristics over time.


\subsection{DyMAB solvers}
To select heuristics dynamically during search, DyMAB employs two non-stationary MAB policies: DyMAB($\epsilon$-Greedy) and DyMAB($\alpha$-UCB). Both policies incorporate the windowed reward estimates and a logarithmically decaying exploration parameter (\(\epsilon\) and \(\alpha\), respectively). This decay schedule gradually shifts the algorithm from exploration to exploitation over time.

\subsubsection{DyMAB($\epsilon$-Greedy) solver}
DyMAB($\epsilon$-Greedy) is presented in Algorithm~\ref{alg:cap1}. Each heuristic $H \in \Omega$ is initialized with an estimated mean reward $\hat{\mu}_H = 0$ and an empty sliding window queue $Q_H$ (line~1). An initial solution $P$ is computed using Prioritized Planning, and the iteration counter $N$ and exploration rate $\epsilon_t$ are initialized (line~2). At each iteration, a heuristic is selected by $\epsilon$-Greedy: with probability $\epsilon_t$ a random heuristic is chosen (exploration), otherwise the heuristic with the highest $\hat{\mu}_H$ is chosen (exploitation) (lines~4--6). The selected heuristic identifies a subset of agents $A_s$; their paths $P^-$ are replanned to produce $P^+$ (lines~7--8), and $P$ is updated if the cost improves (line~9). The reward $r_H = l(P^-) - l(P^+)$ is computed (line~10) and passed to \textsc{UpdateQueue}, which appends it to $Q_H$ and discards the oldest entry if $|Q_H| > W$; $\hat{\mu}_H$ is recomputed as the mean of $Q_H$ (line~11). Finally, $\epsilon_t$ decays logarithmically to shift the policy from exploration toward exploitation over time (line~12).
\begin{algorithm}[t]
\caption{DyMAB($\epsilon$-Greedy)}
\label{alg:cap1}
\KwIn{Graph $G=(V,E)$, agents $A$, heuristic set $\Omega$, $\epsilon$, $\lambda$, $W$}
\KwOut{Collision-free paths $P$}
$\hat{\mu}_H \gets 0$,\ $Q_H \gets [\,]$ \textbf{for each} $H \in \Omega$\;
$P \gets \textsc{PP}(A, G)$;\ $N \gets 0$;\ $\epsilon_t \gets \epsilon$\;
\Repeat{time budget reached}{
  $N \gets N + 1$\;
  \eIf{$\mathrm{rand()} < \epsilon_t$}{
    $H \gets \textsc{RandomHeuristic}(\Omega)$\;
  }{
    $H \gets \arg\max_{H' \in \Omega}\, \hat{\mu}_{H'}$\;
  }
  $A_s \gets H(A)$;\ $P^- \gets \{P_i \mid a_i \in A_s\}$\;
  $P^+ \gets \textsc{Replan}(G, A_s, P \setminus P^-)$\;
  \If{$l(P^+) < l(P^-)$}{$P \gets (P \setminus P^-) \cup P^+$\;}
  $r_H \gets l(P^-) - l(P^+)$\;
  $\textsc{UpdateQueue}(Q_H, r_H, W)$;\ $\hat{\mu}_H \gets \mathrm{mean}(Q_H)$\;
  $\epsilon_t \gets \epsilon / \log(1 + \lambda \cdot N)$\;
}
\Return $P$\;
\end{algorithm}
\noindent\small\textbf{Note:} $\textsc{UpdateQueue}(Q, r, W)$: append $r$ to $Q$; if $|Q| > W$, remove the oldest entry.

\subsubsection{DyMAB($\alpha$-UCB) solver}
DyMAB($\alpha$-UCB) is presented in Algorithm~\ref{alg:cap2}. Each heuristic $H \in \Omega$ is initialized with $\hat{\mu}_H = 0$, a pull count $n_H = 0$, and an empty queue $Q_H$; the exploration coefficient is set to $\alpha_0$ (lines~1--2). An initial solution $P$ is computed using PP (line~2). At each iteration, the heuristic maximizing the UCB index $\hat{\mu}_H + \sqrt{2\alpha/(n_H+1)}$ is selected (line~4). The chosen heuristic identifies $A_s$; paths $P^-$ are replanned to $P^+$ (lines~5--6), and $P$ is updated on improvement (line~7). The reward $r_H = l(P^-) - l(P^+)$ is computed (line~8) and processed by \textsc{UpdateQueue} (shared with Algorithm~\ref{alg:cap1}); $\hat{\mu}_H$ is recomputed and $n_H$ is set to the current queue size $|Q_H|$ (line~9). The coefficient $\alpha$ decays logarithmically each iteration, progressively reducing the exploration bonus (line~10).
\begin{algorithm}[t]
\caption{DyMAB($\alpha$-UCB)}
\label{alg:cap2}
\KwIn{Graph $G=(V,E)$, agents $A$, heuristic set $\Omega$, $\alpha_0$, $\lambda$, $W$}
\KwOut{Collision-free paths $P$}
$\hat{\mu}_H \gets 0$,\ $n_H \gets 0$,\ $Q_H \gets [\,]$ \textbf{for each} $H \in \Omega$\;
$P \gets \textsc{PP}(A, G)$;\ $N \gets 0$;\ $\alpha \gets \alpha_0$\;
\Repeat{time budget reached}{
  $N \gets N + 1$\;
  $H \gets \arg\max_{H' \in \Omega}\!\left(\hat{\mu}_{H'} + \sqrt{\dfrac{2\alpha}{n_{H'}+1}}\right)$\;
  $A_s \gets H(A)$;\ $P^- \gets \{P_i \mid a_i \in A_s\}$\;
  $P^+ \gets \textsc{Replan}(G, A_s, P \setminus P^-)$\;
  \If{$l(P^+) < l(P^-)$}{$P \gets (P \setminus P^-) \cup P^+$\;}
  $r_H \gets l(P^-) - l(P^+)$\;
  $\textsc{UpdateQueue}(Q_H, r_H, W)$;\ $\hat{\mu}_H \gets \mathrm{mean}(Q_H)$;\ $n_H \gets |Q_H|$\;
  $\alpha \gets \alpha_0 / (1 + \lambda \cdot \log(1 + N))$\;
}
\Return $P$\;
\end{algorithm}
%%%%%%%%%%%%%%%%%%%%%%%%%%%%%%%%%%%%%%%%%%%%%%%%%%%%%%%%%%%%%%%%%%%%%%%%

\section{Theoretical analysis of DyMAB}\label{sec:theory}
We now provide a general formal theoretical justification for \text{DyMAB} for a set of $K$ arms or heuristics, showing that under standard assumptions of non-stationary bandits, \text{DyMAB} achieves sub-linear regret, provided that the exploration parameters ($\alpha$ in DyMAB($\alpha$-UCB) and $\epsilon$ in  DyMAB($\epsilon$-Greedy)) decay logarithmically.

\begin{definition}[Non-Stationary MAB for Heuristic Selection]
We consider $H = \{ H_1, H_2, \ldots, H_K \}$ a set of heuristics (arms). At each time $t$, the reward from heuristic $H_i$ is drawn from distribution $D_i(t)$ with unknown and a time-varying mean $\mu_i(t)$.
We assume a bounded reward ($r_t \in [0, 1]$) for simplicity, and piecewise stationary, meaning there is at most $B$ breakpoints dividing the time horizon $T$ into $B \ + \ 1$ stationary segments. \text{DyMAB} use a sliding window of size $W$. 
\end{definition}

\begin{definition}[Dynamic Pseudo-Regret]
The dynamic pseudo-regret incurred by a policy $\pi$ over $T$ rounds (iterations) is defined as:
\begin{equation}
R_T^\pi = \sum_{t=1}^{T} \mu^*(t) - \mathbb{E}[r_t^\pi],
\end{equation}
where $\mu^*(t) = \max_i \mu_i(t)$ is the reward of the best heuristic at time $t$, and $r_t^\pi$ is the reward obtained by policy $\pi$.
\end{definition}

\begin{theorem}[Dynamic Regret Bound of DyMAB($\alpha$-UCB)]\label{thm:dymab_ucb}
Assume that DyMAB($\alpha$-UCB) operates in a piecewise stationary environment with at most $B$ breakpoints and uses a sliding window of size $W$ and an exploration parameter $\alpha$ decaying logarithmically as 
\begin{equation}
    \alpha = \frac{\alpha_0}{(1+ \lambda \log(1 + N_{iteration}))}
\end{equation}
for constants $\alpha_0, \lambda > 0$. Then, the dynamic pseudo-regret of DyMAB($\alpha$-UCB) is bounded as:
\begin{equation}
R_T = \mathcal{O}\left(K \frac{T}{W} \log W + KBW\right).
\end{equation}
\end{theorem}

\begin{proof}
We base our proof directly on Garivier and Moulines' theorem~\cite[Theorem 7]{Non-Stationary}, which states:
\begin{quote}
For any integer $\tau$ and any arm $i \in \{1, \ldots, K\}$:
\begin{equation}
\mathbb{E}_\tau\left[\tilde{N}_T(i)\right] \leq C(\tau) \frac{T \log \tau}{\tau} + \tau \Upsilon_T + \log^2 \tau,
\end{equation}
where $C(\tau)$ is a constant depending on $\tau$ and the reward gaps, and $\Upsilon_T$ is the number of breakpoints. $\tilde{N}_T(i)$ denotes the number of times a sub-optimal arm $i$ has been chosen during the first $T$ round.
\end{quote}
In DyMAB, we use a sliding window of size $W$, corresponding to $\tau$ in the theorem, and we operate in a piecewise stationary environment with at most $B$ breakpoints, thus $\Upsilon_T = B$.

\noindent \text{Step 1: Regret within stationary segments.}

Applying Garivier and Moulines' theorem~\cite[Theorem 7]{Non-Stationary} directly, for any suboptimal heuristic $i$ in any stationary segment:
\begin{equation}
\mathbb{E}\left[\tilde{N}_T(i)\right] \leq \mathcal{O}\left(\frac{T \log W}{W} + BW\right).
\end{equation}

\noindent \text{Step 2: Total regret over all arms.}

Summing over all $K$ arms:
\begin{equation}
R_T = \sum_{i=1}^K \mathbb{E}\left[\tilde{N}_T(i)\right] \leq \mathcal{O}\left(K \frac{T \log W}{W} + KBW\right).
\end{equation}


\noindent \text{Step 3: Conclusion.}

Thus, DyMAB($\alpha$-UCB), using a sliding window of size $W$ achieves:
\begin{equation}
R_T = \mathcal{O}\left(K \frac{T}{W} \log W + KBW\right).
\end{equation}
\end{proof}

In the MAB literature, $K \frac{T}{W} \log W$ correspond to the regret in a stationary segment and $KBW$ to the regret incurred during the breakpoint.


\begin{theorem}[Dynamic Regret Bound of DyMAB($\epsilon$-Greedy)]
Assume that DyMAB($\epsilon$-Greedy) operates in a piecewise stationary environment with at most $B$ breakpoints, using a sliding window of size $W$. Assume an exploration probability $\epsilon_t$ decaying logarithmically as $\epsilon_t = \frac{\epsilon_0}{\log(1 + \lambda t)}$ for constants $\epsilon_0, \lambda > 0$. Then, the dynamic pseudo-regret of DyMAB($\epsilon$-Greedy) is bounded as:
\begin{equation}
R_T = \mathcal{O}\left( \frac{T}{\log T} + K T \cdot \sqrt{\frac{\log W}{W}} + K B W \right).
\end{equation}
\end{theorem}

\begin{proof}
We decompose the total regret into:
\begin{equation}
R_T = R_T^{\text{stationary}} + R_T^{\text{breakpoints}}.
\end{equation}

\noindent \text{Step 1: Regret within Stationary Segments ($R_T^{\text{stationary}}$).}
Consider a stationary segment of length $\tau_s$, during which the reward distributions are fixed.

\noindent \textit{Exploration regret:} \newline
At each round $t$, with probability $\epsilon_t$, DyMAB selects a heuristic uniformly at random. In the worst case, it incurs unit regret per round (i.e., $\Delta_i \leq 1$). Hence, the total exploration regret over the entire horizon $T$ is the sum of regrets incurred at each exploration step:
\begin{equation}
R_T^{\text{explore}} = \sum_{t=1}^{T} \epsilon_t  \leq \sum_{t=1}^{T} \frac{\epsilon_0}{\log(1 + \lambda t)}.
\end{equation}
For large $T$, this sum can be approximated by an integral, and it is known that $\sum_{t=1}^{T} \frac{1}{\log t} = \mathcal{O}\left(\frac{T}{\log T}\right)$.
Therefore, the total exploration regret is:
\begin{equation}
R_T^{\text{explore}} = \mathcal{O}\left(\frac{T}{\log T}\right).
\end{equation}

\noindent \textit{Exploitation regret:} \newline
With probability $1 - \epsilon_t$, DyMAB selects the heuristic with the highest empirical mean over the last $W$ samples. Regret occurs if a suboptimal arm is selected.

Let $i^*$ be the optimal heuristic with mean reward $\mu_{i^*}$, and let $i \neq i^*$ be any suboptimal heuristic with $\mu_i$ and gap $\Delta_i = \mu_{i^*} - \mu_i > 0$. Let $\hat{\mu}_j(W)$ be the empirical mean over the last $W$ samples for heuristic $j$.

We define:
\begin{equation}
\delta_W := \sqrt{\frac{\log W}{2W}}.
\end{equation}

Assume $\Delta_i > 2\delta_W$. If both empirical estimates are close to their means:
\begin{equation}
|\hat{\mu}_i(W) - \mu_i| < \delta_W \quad \text{and} \quad |\hat{\mu}_{i^*}(W) - \mu_{i^*}| < \delta_W,
\end{equation}
then:
\begin{equation}
\hat{\mu}_i(W) < \mu_i + \delta_W = \mu_{i^*} - \Delta_i + \delta_W < \mu_{i^*} - \delta_W < \hat{\mu}_{i^*}(W),
\end{equation}
which contradicts $\hat{\mu}_i(W) \geq \hat{\mu}_{i^*}(W)$, because during exploitation, if the algorithm chooses the suboptimal arm $i$, it means that the empirical mean of arm $i$ looked better than or equal to the empirical mean of the optimal arm $i^*$.

Thus, the event $\hat{\mu}_i(W) \geq \hat{\mu}_{i^*}(W)$ implies that at least one of the two deviations exceeds $\delta_W$:
\begin{equation}
\mathbb{P}\left(\hat{\mu}_i(W) \geq \hat{\mu}_{i^*}(W)\right) \leq \mathbb{P}\left(|\hat{\mu}_i(W) - \mu_i| \geq \delta_W\right) + \mathbb{P}\left(|\hat{\mu}_{i^*}(W) - \mu_{i^*}| \geq \delta_W\right).
\end{equation}

Applying Hoeffding’s inequality (for bounded rewards in $[0,1]$):
\begin{equation}
\mathbb{P}\left(|\hat{\mu}_j(W) - \mu_j| \geq \delta_W\right) \leq 2 \exp(-2W \delta_W^2) = 2 \exp(-\log W) = \frac{2}{W}.
\end{equation}

Therefore:
\begin{equation}
\mathbb{P}\left(\hat{\mu}_i(W) \geq \hat{\mu}_{i^*}(W)\right) \leq \frac{4}{W}.
\end{equation}

\noindent \textit{Expected exploitation regret per round:} \newline
We partition suboptimal arms into two categories:

\begin{itemize}
    \item \text{Large-gap arms} ($\Delta_i > 2\delta_W$):\\
    For these arms, the probability of choosing a sub-optimal one is at most $\frac{4}{W}$ (by Hoeffding's inequality and union bound). Since $\Delta_i \leq 1$, their total contribution to regret per round is:
    \begin{equation}
    \sum_{i : \Delta_i > 2\delta_W} \Delta_i \cdot \frac{4}{W} \leq \frac{4K}{W}.
    \end{equation}

    \item \text{Small-gap arms} ($\Delta_i \leq 2\delta_W$):\\
    These arms are hard to distinguish from the optimal arm, but even if selected, they incur small regret. The total regret per round is bounded by:
    \begin{equation}
    \sum_{i : \Delta_i \leq 2\delta_W} \Delta_i \leq 2K \delta_W = \mathcal{O}\left( K \cdot \sqrt{\frac{\log W}{W}} \right).
    \end{equation}
\end{itemize}

Combining both contributions, the total expected exploitation regret per round is:
\begin{equation}
\mathcal{O}\left( \frac{K}{W} + K \cdot \sqrt{\frac{\log W}{W}} \right).
\end{equation}

Since typically $\sqrt{\frac{\log W}{W}} > \frac{1}{W}$ for practical $W$, we can simplify:
\begin{equation}
\mathcal{O}\left( \frac{K}{W} + K \cdot \sqrt{\frac{\log W}{W}} \right) = \mathcal{O}\left( K \cdot \sqrt{\frac{\log W}{W}} \right).
\end{equation}

\noindent \textit{Total exploitation regret over the entire horizon $T$:} \newline
The expectation of this per-round exploitation regret is taken over $T$ rounds. Since the exploitation probability is $1-\epsilon_t$, and $\epsilon_t$ decays, $1-\epsilon_t \approx 1$ for large $t$. Summing over $T$ rounds, the total exploitation regret is:
\begin{equation}
R_T^{\text{exploit}} = \sum_{t=1}^T (1-\epsilon_t) \cdot \mathcal{O}\left( K \cdot \sqrt{\frac{\log W}{W}} \right) = \mathcal{O}\left( K T \cdot \sqrt{\frac{\log W}{W}} \right).
\end{equation}

\noindent \textit{Total regret in stationary segments ($R_T^{\text{stationary}}$):} \newline
Combining the total exploration and total exploitation regrets:
\begin{equation}
R_T^{\text{stationary}} = R_T^{\text{explore}} + R_T^{\text{exploit}} = \mathcal{O}\left( \frac{T}{\log T} + K T \cdot \sqrt{\frac{\log W}{W}} \right).
\end{equation}

\noindent \text{Step 2: Regret Due to Breakpoints ($R_T^{\text{breakpoints}}$).}

After each of the $B$ breakpoints, each heuristic may continue to rely on outdated samples for up to $W$ rounds. In the worst case, this induces regret of up to $1$ per round per heuristic. Summing over all $B$ breakpoints and $K$ arms:
\begin{equation}
R_T^{\text{breakpoints}} = \mathcal{O}(K B W).
\end{equation}

\noindent \text{Step 3: Total Regret.}

Combining both the stationary and breakpoint components:
\begin{equation}
R_T = R_T^{\text{stationary}} + R_T^{\text{breakpoints}} = \mathcal{O}\left( \frac{T}{\log T} + K T \cdot \sqrt{\frac{\log W}{W}} + K B W \right).
\end{equation}
\end{proof}

\begin{proposition}[Optimal Window for DyMAB($\alpha$-UCB)]
If the time horizon $T$ and the number of breakpoints $B$ in a piecewise-stationary environment, are known in advance, the window size $W$ can be chosen so as to minimize regret bound for DyMAB($\alpha$-UCB). 
\begin{equation}
W^\star = \Theta\left( \sqrt{ \frac{T \log T}{B} } \right).
\end{equation}
Substituting this value yields:
\begin{equation}
\mathcal{R}_T^{\mathrm{UCB}} = \mathcal{O}\left( K \sqrt{T B \log T} \right).
\end{equation}
\end{proposition}

\vspace{1em}

\begin{proposition}[Optimal Window for DyMAB($\epsilon$-Greedy)]
If the time horizon $T$ and the number of breakpoints $B$ in a piecewise-stationary environment, are known in advance, the window size $W$ can be chosen so as to minimize regret bound for DyMAB($\epsilon$-Greedy).
\begin{equation}
W^\star = \Theta\left( \left( \frac{T^2 \log T}{B^2} \right)^{1/3} \right).
\end{equation}
This yields the regret:
\begin{equation}
\mathcal{R}_T^{\epsilon\text{-Greedy}} = \mathcal{O}\left( \frac{T}{\log T} + K T^{2/3} (\log T)^{1/3} B^{1/3} \right).
\end{equation}
\end{proposition}

%\vspace{1em}

\begin{table}[t]
\centering
\caption{Comparison of optimal regret and window size for DyMAB policies}
\begin{tabular}{|c|c|c|}
\hline
\textbf{Policy} & \textbf{Regret Bound} & \textbf{Optimal Window Size $W^\star$} \\
\hline
DyMAB($\alpha$-UCB) & $\mathcal{O}\left( K \sqrt{T B \log T} \right)$ & $\Theta\left( \sqrt{ \frac{T \log T}{B} } \right)$ \\
\hline
DyMAB($\epsilon$-Greedy) & $\mathcal{O}\left( \frac{T}{\log T} + K T^{2/3} (\log T)^{1/3} B^{1/3} \right)$ & $\Theta\left( \left( \frac{T^2 \log T}{B^2} \right)^{1/3} \right)$ \\
\hline
\end{tabular}
\label{tab:comp}
\end{table}

As presented in Table~\ref{tab:comp}, the regret bounds for the two DyMAB policies differ in their asymptotic behavior. Given an optimal window size $W^*$, the regret for DyMAB($\alpha$-UCB) can be minimized up to  $\mathcal{O}( K \sqrt{T B \log T})$, which is a well known bound for sliding window UCB \citep{Non-Stationary, Non-Stationary3}. For DyMAB($\epsilon$-Greedy), the regret bound consists of two terms: $\mathcal{O} \left( \frac{T}{\log T} + K T^{2/3} (\log T)^{1/3} B^{1/3} \right),$
where the first term accounts for the cumulative regret due to uniform random exploration, and the second term reflects the exploitation regret under non-stationarity.

Both policies achieve sublinear regret in $T$, but differ in their sensitivity to the number of breakpoints $B$ and in the structure of their exploration-exploitation trade-offs. DyMAB($\alpha$-UCB) exhibits a tighter dependence on both $T$ and $B$—scaling as $\sqrt{T B}$, which makes it theoretically more favorable in environments with infrequent distributional shifts. In contrast, DyMAB($\epsilon$-Greedy)'s regret grows more rapidly with $T$, but its structure is analytically simpler and does not require confidence-based computations. Therefore, from a purely theoretical standpoint, DyMAB($\alpha$-UCB) offers a stronger regret guarantee when the number of breakpoints $B$ and the horizon $T$ are known.

%%%%%%%%%%%%%%%%%%%%%%%%%%%%%%%%%%%%%%%%%%%%%%%%%%%%%%%%%%%%%%%%%%%%%%%%
\section{Experimental evaluation} \label{sec:experiments}
We evaluate $\text{DyMAB}$ algorithms on several maps from the Moving AI Lab benchmarks \citep{SternEtal2019}, as well as our newly introduced city map, $\texttt{NewCity(|V|=572000)}$\footnote{github repository link redacted for review process}. Compared to the standard benchmark maps, the NewCity\footnote{Linköping city map} map is significantly larger, increasing the robustness of our evaluation.
In all experiments, DyMAB uses exactly three heuristics: $H_{\text{rand}}$, $H_{\text{sync}}$, and $H_{\text{entr}}$. The original MAPF-LNS~\citep{LNS} heuristics $H_{\text{agent}}$ and $H_{\text{grid}}$ are deliberately excluded. Following the ALNS orthogonality principle~\citep{LNS2}, the pool should consist of operators covering distinct search dimensions: $H_{\text{sync}}$ subsumes $H_{\text{agent}}$ (both target delayed agents), and $H_{\text{entr}}$ subsumes $H_{\text{grid}}$ (both target spatially congested regions). Including all five would introduce redundancy without additional complementary coverage.
To implement $\alpha$-UCB$(\alpha, \lambda_1, W_1)$ and $\epsilon$-Greedy$(\epsilon, \lambda_2, W_2, T)$, we performed a grid search to tune all hyperparameters. Unless stated otherwise, we use the following default values: exploration parameters $\alpha=10000$ and $\epsilon=0.5$, decay factors $\lambda_1=10$ and $\lambda_2=0.01$, reward queue sizes $W_1=W_2=100$, and an exploration period $T=1000$. To evaluate $\text{DyMAB}$, we conduct experiments across 25 random scenarios for each map, and report the average Sum of Delays and AUCs. All experiments were run on an Intel Core i9-10900X CPU with 128GB RAM. Implementations of the proposed algorithms are available as an open-source project [xx].

\begin{figure}[t]
    \centering
    \includegraphics[width=0.9\textwidth]{./images/experiment1.png}
    \caption{Average Sum of Delays for DyMAB($\alpha$-UCB) and DyMAB($\epsilon$-Greedy) with different time budgets compared with the empirical best Policy. The error bar shows the $95\%$ confidence interval. The runtime represents the time budget in seconds.}
    \label{fig:exp1}
\end{figure}

\subsection{Experiment 1: DyMAB convergence}
This experiment aims to assess the convergence properties of $\text{DyMAB}$ when the time budget increases. To achieve this, we plot the Sum of Delays with respect to the time budget on $\texttt{Dense520d (|V| = 28178)}$ and \texttt{Berlin\_1\_256 (|V| = 47540)} maps, with 400 and 600 agents respectively. We compared the results with the empirical best policy which is determined as the minimum Sum of Delays across all configurations. As shown in Fig~\ref{fig:exp1}, all DyMAB variants converge towards the near-optimal empirical best policy as the time budget (runtime) increases. Both DyMAB($\alpha$-UCB) and DyMAB($\epsilon$-Greedy) demonstrate comparable performance across the tested maps.                                                                                                                                                                                                                                                                                                                                                                                                                                                                                                                                                                                                                                                                                                                                                                                                                                                                                                                                                                                                                                                                                                                                                                                                                                                                                                                                                                                                                                                                                                                                                                                                                                                                                                                                                                                                                                                                                                                                                                                                                                                                                                                                                                                                                                                                                                                                                                                                                                                                                                                                                                                                                                                                                                                                                                                                                                                                                                                                                                                                                                                                          

\begin{table*}[t]
\renewcommand{\arraystretch}{1.4}
\centering
\setlength{\tabcolsep}{1pt}
\caption{Average Sum of Delays and AUCs for different heuristics: Random Neighbourhood ($H_{rand}$), Grid-Based Neighbourhood ($H_{grid}$), Agent-Based Neighbourhood ($H_{agent}$), Synchronised RandomWalk Neighbourhood ($H_{sync}$), Entropy-Based Neighbourhood ($H_{entr}$), and Adaptive using DyMAB($\alpha$-UCB) policy. The time budget is fixed at 100 seconds. Bold indicates the best value for Cost and AUC, and italic indicates the second-best.}
\resizebox{\textwidth}{!}{%
\begin{tabular}{lccc|cc|cc|cc|cc|cc|cc|cc}
    \toprule
    \textbf{Map} & \textbf{m} & \multicolumn{2}{c|}{\textbf{$H_{rand}$}} & \multicolumn{2}{c|}{\textbf{$H_{grid}$}} & \multicolumn{2}{c|}{\textbf{$H_{agent}$}} & \multicolumn{2}{c|}{\textbf{$H_{sync}$}} & \multicolumn{2}{c|}{\textbf{$H_{entr}$}} & \multicolumn{2}{c|}{\textbf{DyMAB($\alpha$-UCB)}} \\
    \cmidrule(lr){3-4} \cmidrule(lr){5-6} \cmidrule(lr){7-8} \cmidrule(lr){9-10} \cmidrule(lr){11-12} \cmidrule(lr){13-14}
    &  & \text{Cost} & \text{AUC} & \text{Cost} & \text{AUC} & \text{Cost} & \text{AUC} & \text{Cost} & \text{AUC} & \text{Cost} & \text{AUC} & \text{Cost} & \text{AUC} \\
    \midrule
    \multirow{3}{*}{\rotatebox{90}{Random}} & 100 & 4339.08 & \textit{931} & 4330.76 & 933 & 4331.52 & 957 & 4331.20 & 953 & \textbf{4313.24} & 2412 & \textit{4330.48} & \textbf{846} \\
    & 200 & 8674.68 & 6685 & 8553.64 & 5806 & 8554.84 & 4226 & \textit{8552.88} & \textit{4011} & 8615.04 & 9154 & \textbf{8550.60} & \textbf{3834} \\
    & 300 & 13003.29 & 37664 & 13016.71 & 22922 & 12899.04 & 12762 & \textit{12892.75} & \textit{12454} & 13002.41 & 24832 & \textbf{12891.71} & \textbf{12431} \\
    \hline
    \multirow{3}{*}{\rotatebox{90}{Dense}} & 200 & 34705.56 & 46379 & 34875.72 & 51775 & \textit{34616.24} & \textbf{7444} & 34640.72 & 15115 & 34765.16 & 28792 & \textbf{34604.36} & \textit{9082} \\
    & 400 & 72303.76 & 451169 & 71374.24 & 329440 & \textit{69491.40} & \textbf{93417} & \textbf{69490.56} & 114784 & 70493.92 & 261053 & 69597.08 & \textit{102949} \\
    & 600 & 114151.08 & 1239469 & 110460.48 & 995797 & \textbf{105949.52} & 501808 & 106262.80 & 552941 & 110219.32 & 981363 & \textit{105992.76} & \textbf{455550} \\
    \hline
    \multirow{3}{*}{\rotatebox{90}{Berlin}} & 500 & 94269.08 & 540660 & 92313.88 & 339880 & 90289.52 & \textit{75805} & \textit{90287.96} & 80153 & 91478.76 & 295241 & \textbf{90250.52} & \textbf{52460} \\
    & 800 & 158039.40 & 1544684 & 153406.36 & 1186468 & 147368.88 & 660376 & \textit{147360.20} & \textit{660087} & 153806.24 & 1259103 & \textbf{146125.52} & \textbf{442094} \\
    & 1000 & 202849.91 & 2390499 & 201648.13 & 2249260 & 190882.13 & 1453086 & \textit{190456.25} & \textit{1499752} & 199148.04 & 2158026 & \textbf{186354.17} & \textbf{1082926} \\
    \bottomrule
\end{tabular}
}
\label{tab:performance_metrics}
\end{table*}

\begin{comment}
    

\renewcommand{\arraystretch}{1.4}
\begin{table*}[t]
    \centering
    \setlength{\tabcolsep}{0.8pt}
        \caption{Average Sum of Delays and AUCs for different heuristics: Random Neighbourhood ($H_{rand}$), Grid-Based Neighbourhood ($H_{grid}$), Agent-Based Neighbourhood ($H_{agent}$), Synchronised RandomWalk Neighbourhood ($H_{sync}$), Entropy-Based Neighbourhood ($H_{entr}$), and Adaptive using DyMAB($\alpha$-UCB) policy. The time budget is fixed at 100 seconds. Bold indicates the best value for Cost and AUC, and italic indicates the second-best.}
    \begin{tabular}{lccc|cc|cc|cc|cc|cc|cc|cc}
        \toprule
        \textbf{Maps} & \textbf{m} & \multicolumn{2}{c|}{\textbf{$H_{rand}$}} & \multicolumn{2}{c|}{\textbf{$H_{grid}$}} & \multicolumn{2}{c|}{\textbf{$H_{agent}$}} & \multicolumn{2}{c|}{\textbf{$H_{sync}$}} & \multicolumn{2}{c|}{\textbf{$H_{entr}$}} & \multicolumn{2}{c|}{\textbf{DyMAB($\alpha$-UCB)}} \\
        \cmidrule(lr){3-4} \cmidrule(lr){5-6} \cmidrule(lr){7-8} \cmidrule(lr){9-10} \cmidrule(lr){11-12} \cmidrule(lr){13-14}
        &  & \text{Cost} & \text{AUC} & \text{Cost} & \text{AUC} & \text{Cost} & \text{AUC} & \text{Cost} & \text{AUC} & \text{Cost} & \text{AUC} & \text{Cost} & \text{AUC} \\
        \midrule
        \multirow{3}{*}{\rotatebox{90}{Random}} & 100 & 4339.08 & \textit{931} & 4330.76 & 933 & 4331.52 & 957 & 4331.20 & 953 & \textbf{4313.24} & 2412 & \textit{4330.48} & \textbf{846} \\
        & 200 & 8674.68 & 6685 & 8553.64 & 5806 & 8554.84 & 4226 & \textit{8552.88} & \textit{4011} & 8615.04 & 9154 & \textbf{8550.60} & \textbf{3834} \\
        & 300 & 13003.29 & 37664 & 13016.71 & 22922 & 12899.04 & 12762 & \textit{12892.75} & \textit{12454} & 13002.41 & 24832 & \textbf{12891.71} & \textbf{12431} \\
        \hline
        \multirow{3}{*}{\rotatebox{90}{Dense}} & 200 & 34705.56 & 46379 & 34875.72 & 51775 & \textit{34616.24} & \textbf{7444} & 34640.72 & 15115 & 34765.16 & 28792 & \textbf{34604.36} & \textit{9082} \\
        & 400 & 72303.76 & 451169 & 71374.24 & 329440 & \textit{69491.40} & \textbf{93417} & \textbf{69490.56} & 114784 & 70493.92 & 261053 & 69597.08 & \textit{102949} \\
        & 600 & 114151.08 & 1239469 & 110460.48 & 995797 & \textbf{105949.52} & 501808 & 106262.80 & 552941 & 110219.32 & 981363 & \textit{105992.76} & \textbf{455550} \\
        \hline
        \multirow{3}{*}{\rotatebox{90}{Berlin}} & 500 & 94269.08 & 540660 & 92313.88 & 339880 & 90289.52 & \textit{75805} & \textit{90287.96} & 80153 & 91478.76 & 295241 & \textbf{90250.52} & \textbf{52460} \\
        & 800 & 158039.40 & 1544684 & 153406.36 & 1186468 & 147368.88 & 660376 & \textit{147360.20} & \textit{660087} & 153806.24 & 1259103 & \textbf{146125.52} & \textbf{442094} \\
        & 1000 & 202849.91 & 2390499 & 201648.13 & 2249260 & 190882.13 & 1453086 & \textit{190456.25} & \textit{1499752} & 199148.04 & 2158026 & \textbf{186354.17} & \textbf{1082926} \\
        \bottomrule
    \end{tabular}
    \label{tab:performance_metrics}
\end{table*}
\end{comment}
\subsection{Experiment 2: Neighborhood selection}
Here, we compare the performance of the three heuristics discussed in Section 4, such as $H_{sync},\ H_{entr} \ \text{and} \ H_{rand}$ against Agent-Based Neighborhood ($H_{agent}$), and Grid-Based Neighborhood ($H_{grid}$) presented in \text{MAPF-LNS} framework \citep{LNS}. The experiments are conducted on $\texttt{random-64-64-20 (|V| = 3270)}$ map and the two maps used in the previous experiment. The results, detailed in Table~\ref{tab:performance_metrics}, indicate that the performance of the heuristics varies depending on the specific map and the number of agents $m$. Among the heuristics, $H_{sync}$ demonstrates consistent performance, achieving lower Cost and AUC values compared to $H_{rand}$, $H_{agent}$, and $H_{grid}$, particularly in dense and large-scale scenarios. $H_{entr}$ frequently yielded the lowest Cost and AUC, leveraging information-theoretic measures to identify critical paths and agents.
The Adaptive version, corresponding to DyMAB($\alpha$-UCB), demonstrates the best overall performance, maintaining consistently low Cost and AUC. This is expected, as DyMAB learns to select the most effective heuristic over time. If we choose DyMAB($\epsilon$-Greedy), similar performance will be observed.

\subsection{Experiment 3: Effect of Sliding Window Size $W$} \label{sec:window}

This experiment evaluates the impact of the sliding window size $W$ on the performance of DyMAB($\alpha$-UCB) and DyMAB($\epsilon$-Greedy). The windowed reward queue $Q$ plays a central role in enabling DyMAB to adapt to non-stationary environments by retaining recent reward observations for each heuristic. A smaller window size enables faster adaptation to changes in heuristic performance by emphasizing recent data, while a larger window provides more stable estimates, at the cost of slower responsiveness to changes.

All experiments are conducted with a time budget of 120 seconds on the Dense and Berlin maps, using varying numbers of agents. The average Sum of Delays over 25 random scenarios is plotted in Fig.~\ref{fig:combined}.

\begin{figure}[t]
    \centering
    %\begin{subfigure}[t]{0.48\textwidth}
    \begin{subfigure}[t]{0.7\textwidth}
        \centering 
        \includegraphics[width=\textwidth]{./images/experiment22.png}
        \caption{}
        \label{fig:exp22}
    \end{subfigure}
    \hfill
    \begin{subfigure}[t]{0.7\textwidth}
        \centering
        \includegraphics[width=\textwidth]{./images/experiment21.png}
        \caption{}
        \label{fig:exp21}
    \end{subfigure}
    %\vspace{1em}
    \caption{Average Sum of Delays for DyMAB($\alpha$-UCB) and DyMAB($\epsilon$-Greedy) with different sliding window sizes $W$, and respective exploration parameters $\alpha$ and $\epsilon$. (a) Berlin map. (b) Dense map.}
    \label{fig:combined}
\end{figure}

DyMAB($\alpha$-UCB) benefits from larger window sizes, particularly when paired with higher exploration parameters (e.g., $\alpha = 50000$ or $\alpha = 100000$). Larger windows yield more reliable empirical mean estimates and tighter confidence bounds, which improve decision-making in environments where heuristic effectiveness changes frequently. This is consistent with the theoretical regret bound $\mathcal{O}(K \sqrt{T B \log T})$, which is minimized when $W = \Theta(\sqrt{T \log T / B})$. In such cases, higher $\alpha$ values promote sufficient exploration early on, allowing the algorithm to identify the optimal heuristics with greater confidence.
In contrast, DyMAB($\epsilon$-Greedy) exhibits stable performance across different values of $W$ and $\epsilon$, indicating greater robustness to parameter tuning. The random exploration ensures all heuristics are periodically sampled, making the method less sensitive to the precision of reward estimates. These results suggest that while DyMAB($\alpha$-UCB) is more sensitive to the proper tuning of $W$ and $\alpha$, it may achieve better performance when such tuning is feasible. Conversely, DyMAB($\epsilon$-Greedy) offers a more robust alternative under less favorable tuning conditions.

\subsection{Experiment 4: Neighborhood Size M}
We evaluate the performance of DyMAB and three heuristic methods under varying neighborhood sizes $M$ across different maps and number of agents $m$. Figure~\ref{fig:exp5} illustrates the average Sum of Delays across 25 scenarios for each map. As shown, the performance of the heuristics improves with increasing neighborhood size $M$. However, DyMAB demonstrates consistent performance regardless of $M$, achieving the smallest Sum of Delays across all scenarios. This stability is attributed to DyMAB’s ability to focus on the most effective heuristic dynamically.
\begin{figure*}[t]
    \centering
    \includegraphics[width=0.7\textwidth]{./images/experiment5.png}
    \vspace{1em}
    \caption{Sum of delays for Random Neighborhood (RN), Synchronized RandomWalk Neighborhood (SRWN), Entropy-Based Neighborhood (EBN),
and Adaptive using DyMAB($\alpha$-UCB) for different neighborhood sizes. Shaded areas show the $95\%$ confidence interval. The legend applies to all plots.}
    \label{fig:exp5}
\end{figure*}

\begin{figure*}[t]
    \centering
    \includegraphics[width=\textwidth]{./images/experiment3.png}
    \caption{Sum of delays for DyMAB($\alpha$-UCB) and DyMAB($\epsilon$-Greedy) compared with BALANCE (Thompson), BALANCE (UCB1), MAPF-LNS and MAPF-LNS2 for different numbers of agents. All experiments are run on the same hardware configuration with a time budget of 120 seconds. Shaded areas show the $95\%$ confidence interval. The legend applies to all plots.}
    \label{fig:exp3}
\end{figure*}

\begin{table*}[t]
    \centering
    \setlength{\tabcolsep}{0.7pt}
    \caption{Average iterations and AUCs for all algorithms. "Iter" refers to the number of iterations. Bold indicates the best value for AUC.}
    
    \begin{tabular}{l|c|cc|cc|cc|cc|cc|cc|cc}
        \toprule
        \textbf{Map} & \textbf{m} & \multicolumn{2}{c|}{\textbf{MAPF-LNS}} & \multicolumn{2}{c|}{\textbf{MAPF-LNS2}} & \multicolumn{2}{c|}{\textbf{Thompson}} & \multicolumn{2}{c|}{\textbf{UCB1}} & \multicolumn{2}{c|}{\textbf{DyMAB($\alpha$-UCB)}} & \multicolumn{2}{c}{\textbf{DyMAB($\epsilon$-Greedy)}} \\
        \cmidrule(lr){3-4} \cmidrule(lr){5-6} \cmidrule(lr){7-8} \cmidrule(lr){9-10} \cmidrule(lr){11-12} \cmidrule(lr){13-14}
        &  & \text{Iter} & \text{AUC} & \text{Iter} & \text{AUC} & \text{Iter} & \text{AUC} & \text{Iter} & \text{AUC} & \text{Iter} & \text{AUC} & \text{Iter} & \text{AUC} \\
        \midrule
        \multirow{3}{*}{\rotatebox{90}{Random}} & 100 & 16384 & 1134 & 154801 & 1027 & 192583 & 1705 & 185417 & 1638 & 62079 & 1021 & 61369 & \textbf{1014} \\
        & 200 & 4250 & 12054 & 70176 & 5059 & 82540 & 19550 & 79658 & 18928 & 22128 & 4559 & 21937 & \textbf{4535} \\
        & 300 & 1964 & 95324 & 38010 & 18925 & 42659 & 116424 & 41248 & 112617 & 13474 & 14475 & 13454 & \textbf{14260} \\
        \midrule
        \multirow{3}{*}{\rotatebox{90}{Ost003d}} & 200 & 297 & 269581 & 11839 & 40477 & 17576 & 140851 & 18106 & 134394 & 2282 & 51300 & 2306 & \textbf{50964} \\
        & 400 & 87 & 1408665 & 5241 & 505958 & 6544 & 1062745 & 6754 & 1012631 & 1249 & 430504 & 1253 & \textbf{429507} \\
        & 600 & 35 & 2841247 & 2479 & 1887567 & 2907 & 2826113 & 2983 & 2706886 & 635 & \textbf{1666151} & 628 & 1671842 \\
        \midrule
        \multirow{3}{*}{\rotatebox{90}{Dense}} & 200 & 105 & 166418 & 12501 & 10766 & 16700 & 46760 & 16863 & 44696 & 3223 & \textbf{9818} & 3061 & 10125 \\
        & 500 & 24 & 1397331 & 3968 & 405961 & 4547 & 932364 & 4586 & 892128 & 1144 & \textbf{246785} & 1048 & 261397 \\
        & 800 & 11 & 3050654 & 1829 & 2019158 & 1874 & 2820094 & 1901 & 2701834 & 736 & \textbf{1263804} & 671 & 1315039 \\
        \midrule
        \multirow{4}{*}{\rotatebox{90}{Berlin}} & 200 & 84 & 80823 & 16029 & \textbf{3528} & 19857 & 19330 & 18790 & 18625 & 5913 & 6451 & 6007 & 5475 \\
        & 500 & 17 & 897086 & 4768 & 124255 & 5184 & 439184 & 4888 & 428000 & 1613 & 56956 & 1641 & \textbf{56260} \\
        & 800 & 9 & 2089424 & 2125 & 994106 & 2125 & 1724541 & 2003 & 1663081 & 901 & 479939 & 917 & \textbf{474174} \\
        & 1000 & 7 & 3163219 & 1373 & 2027183 & 1274 & 2744739 & 1224 & 2633097 & 671 & 1205120 & 685 & \textbf{1194283} \\
        \midrule
        \multirow{5}{*}{\rotatebox{90}{NewCity}} & 100 & 34 & 40711 & 2791 & \textbf{1195} & 3643 & 6631 & 3477 & 6366 & 2031 & 19846 & 1775 & 14580 \\
        & 500 & 4 & 1252670 & 846 & \textbf{169996} & 1001 & 350159 & 932 & 346250 & 270 & 328572 & 279 & 324671 \\
        & 1000 & 2 & 3717274 & 333 & 2225756 & 363 & 3089092 & 344 & 2939158 & 113 & \textbf{2038848} & 113 & 2080958 \\
        & 1500 & 2 & 6807155 & 155 & 5759002 & 160 & 6687516 & 150 & 6220139 & 66 & \textbf{5445558} & 64 & 5461020 \\
        & 2000 & 2 & 8891196 & 71 & 8987055 & 64 & 8626836 & 61 & 8079584 & 38 & 8979265 & 35 & \textbf{8532032} \\
        \bottomrule
    \end{tabular}
    \label{tab:comparison}
\end{table*}
\subsection{Experiment 5: State-of-the-art comparison}
This experiment compares \text{DyMAB} against the original \text{MAPF-LNS}, \text{MAPF-LNS2}, known to be very efficient, and two stationary MAB policies from BALANCE: \text{Thompson Sampling} and \text{UCB1}. 
We run all algorithms on five different maps such as \texttt{random-64-64-20(|V|=3270)}, \texttt{ostd003d(|V|=13214)}, \texttt{Dense520d(|V|=28178)}, \texttt{Berlin\_1\_256(|V|=47540)} and \texttt{NewCity(|V|=572000)} with different numbers of agents $m$ and a time budget of 120 seconds. 
The results presented in Fig.~\ref{fig:exp3} show that all DyMAB policies (DyMAB($\alpha$-UCB) and DyMAB($\epsilon$-Greedy))  outperformed all other algorithms, especially in scenarios with a large number of agents. MAPF-LNS2 performs reasonably well in all maps compared to BALANCE and the original MAPF-LNS. 
DyMAB demonstrates adaptability, dynamically adjusting its heuristic selection through the sliding window rewards mechanism to respond effectively to changes in heuristic performance. Furthermore, DyMAB’s convergence relies on its capability to learn to choose the best heuristic and a logarithmic decay that reduces the exploration over time. Additionally, in complex scenarios with larger maps and increasing agent counts, DyMAB achieves significant reductions in the Sum of Delays, outperforming recent state-of-the-art solvers in both computational efficiency and solution quality. 

In Table~\ref{tab:comparison}, we report the number of iterations and AUCs for each algorithm. DyMAB consistently delivers the smallest AUC values, highlighting its rapid convergence across all tested maps and agent configurations. DyMAB($\epsilon$-Greedy) performs comparably to DyMAB($\alpha$-UCB), often within $1-2\%$ of AUC differences.
DyMAB demonstrates adaptability, dynamically adjusting its heuristic selection through the sliding window rewards mechanism to respond effectively to changes in heuristic performance. Furthermore, DyMAB’s convergence relies on his capability to learn to choose the best heuristic and a logarithmic decay that reduces the exploration over time. Additionally, in complex scenarios with larger maps and increasing agent counts, DyMAB achieves significant reductions in the Sum of Delays and AUC, outperforming recent state-of-the-art solvers in both computational efficiency and solution quality. 

\begin{table*}[t]
    \centering
    \setlength{\tabcolsep}{2pt}
    \caption{Solution Cost and AUC with different large-scale maps ($m=3500$ agents, 500-second time budget). All values are divided by 1000. Bold indicates the best value for Solution Cost and AUC; italic indicates the second-best. LaCAM~\citep{LaCam} and LaCAM*~\citep{LaCamStar} are complete search-based solvers; AUC is not applicable ($-$).}
    \resizebox{\textwidth}{!}{%
    \begin{tabular}{l|cc|cc|cc|cc|cc|c|c|cc}
        \toprule
        \textbf{Map} & \multicolumn{2}{c|}{\textbf{MAPF-LNS2}} & \multicolumn{2}{c|}{\textbf{UCB1}} & \multicolumn{2}{c|}{\textbf{Thompson}} & \multicolumn{2}{c|}{\textbf{DyMAB($\alpha$-UCB)}} & \multicolumn{2}{c|}{\textbf{DyMAB($\epsilon$-Greedy)}} & \textbf{LaCAM} & \textbf{LaCAM*} & \multicolumn{2}{c}{\textbf{ADDRESS}} \\
        \cmidrule(lr){2-3} \cmidrule(lr){4-5} \cmidrule(lr){6-7} \cmidrule(lr){8-9} \cmidrule(lr){10-11} \cmidrule(lr){12-13} \cmidrule(lr){14-15}
        & Cost & AUC & Cost & AUC & Cost & AUC & Cost & AUC & Cost & AUC & Cost & Cost & Cost & AUC \\
        \midrule
        Warehouse & 942.10 & 147109 & 942.91 & 147357 & 942.91 & 147343 & \textbf{931.31} & \textbf{140329} & \textit{937.99} & 142100 & 1069.89 & 1069.89 & 938.36 & \textit{141082} \\
        Boston    & \textit{3453.63} & 152310 & 3467.10 & 152834 & 3467.10 & 152836 & \textbf{3450.02} & \textit{127233} & 3460.53 & 128139 & 3659.15 & 3658.54 & 3457.87 & \textbf{125487} \\
        Moscow    & 3570.87 & 168560 & 3582.12 & 177995 & 3582.12 & 178942 & \textbf{3555.74} & \textbf{146424} & 3578.85 & 151337 & 3796.03 & 3797.79 & \textit{3562.48} & \textit{151229} \\
        New York  & 3527.64 & 116448 & 3529.80 & 118252 & 3529.80 & 118754 & \textbf{3513.80} & \textbf{73414} & 3527.61 & \textit{84457} & 3739.48 & 3734.40 & \textit{3517.13} & 84699 \\
        Paris     & 3418.99 & 148593 & 3432.21 & 141770 & 3432.18 & 156264 & \textbf{3406.73} & \textbf{127768} & 3431.01 & \textit{130885} & 3616.78 & 3616.78 & \textit{3415.79} & 132372 \\
        Shanghai  & 3628.19 & 113580 & 3636.51 & 143298 & 3637.65 & 138719 & \textit{3621.27} & \textbf{99255} & 3638.58 & 100647 & 3877.59 & 3868.83 & \textbf{3620.23} & \textit{100477} \\
        Sydney    & \textbf{3575.25} & 159462 & 3595.47 & 174229 & 3595.47 & 175640 & \textit{3583.26} & \textbf{142397} & 3587.79 & \textit{144329} & 3822.20 & 3821.29 & 3587.48 & 144438 \\
        \bottomrule
    \end{tabular}
    }
    \label{tab:comparison1}
\end{table*}

\subsection{Experiment 6: Large-scale test}
In this experiment, we evaluate \text{DyMAB} on complex city maps and a large warehouse map to assess its performance in large-scale scenarios. The tested maps include 
\texttt{Boston\_0\_1024 (|V| = 796896)}, 
\texttt{Moscow\_0\_1024 (|V| = 795559)}, 
\texttt{NewYork\_0\_1024 (|V| = 792676)}, 
\texttt{Paris\_1\_1024 (|V| = 800729)}, 
\texttt{Shanghai\_0\_1024 (|V| = 789333)}, 
\texttt{Sydney\_0\_1024 (|V| = 793424)}, and 
\texttt{warehouse-20-40-10-2-2 (|V| = 38756)} maps. 
 The number of agents is fixed at $m=3500$. 
Table~\ref{tab:comparison1} presents the results, comparing DyMAB with \text{MAPF-LNS2}, \text{UCB1}, Thompson Sampling, LaCAM~\citep{LaCam}, LaCAM*~\citep{LaCamStar} (both complete search-based solvers), and ADDRESS~\citep{ADDRESS} (an anytime LNS solver using Thompson Sampling with a diverse heuristic pool).
DyMAB($\alpha$-UCB) consistently achieves competitive solution costs across all maps, and demonstrates strong anytime performance reflected by its low AUC values. ADDRESS, which uses Thompson Sampling with a diverse heuristic pool, achieves faster early convergence on some maps (lower AUC on Boston) and a slightly lower final cost on Shanghai, while DyMAB's non-stationary adaptive selection yields better or comparable final solutions overall. LaCAM and LaCAM* serve as reference points for complete search-based approaches; their higher costs confirm the advantage of anytime LNS methods at this scale. LaCAM* produces results similar to LaCAM at this scale because the initial solution search exhausts the majority of the time budget, leaving little room for refinement.

%%%%%%%%%%%%%%%%%%%%%%%%%%%%%%%%%%%%%%%%%%%%%%%%%%%%%%%%%%%%%%%%%%%%%%%%
\section{Conclusion} \label{sec:conclusion}
We have presented DyMAB as an anytime MAPF framework that addresses the non-stationary nature of heuristic performance for ALNS-based solver. We proposed two heuristics, $H_{sync}$ and $H_{entr}$, that explore the solution space efficiently and two policies DyMAB($\alpha$-UCB) and DyMAB($\epsilon$-Greedy), that dynamically learn to choose the best heuristic while taking account of the heuristic shift in performance over time. This adaptability, coupled with their exploration-exploitation strategies(UCB with $\alpha$ decaying and $\epsilon$-Greedy with $\epsilon$ decaying), ensures that DyMAB consistently produces near-optimal solutions within constrained time budgets.
Our theoretical analysis demonstrates that DyMAB achieves sublinear regret under standard non-stationarity assumptions, offering formal justification for its effectiveness for ALNS-based MAPF solvers. 

Our empirical evaluation demonstrates that \text{DyMAB} consistently outperforms recent state-of-the-art MAPF solvers, including \text{MAPF-LNS}, \text{MAPF-LNS2}, \text{LaCAM}, and \text{BALANCE} (using both Thompson Sampling and UCB1), across a diverse range of benchmark scenarios. It is also comparatively most effective at large-scale scenarios, presenting the lowest Sum of Delays and AUC. 

Although \text{DyMAB} framework is tailored to MAPF, its underlying methodology is generalizable to other combinatorial domains involving concurrent heuristic selection, such as Vehicle Routing Problems(VRP), Multi-Agent Pickup and Delivery (MAPD).
%

Furthermore, even if \text{DyMAB} effectively selects among a given set of heuristics, its overall performance is still fundamentally dependent on the quality and diversity of the provided neighborhood heuristics ($H_{sync}$, $H_{entr}$, $H_{rand}$). The development of even more effective heuristics could further enhance the performance. The current \text{DyMAB} framework focuses on non-stationarity within the search process for a static MAPF instance. A natural extension would be to adapt DyMAB to handle fully dynamic MAPF problems, where the environment itself changes over time (e.g., new obstacles, changing agent goals). From a modeling perspective, another promising direction is to consider richer formulations for heuristic selection: a full MDP formulation would explicitly account for the policy-dependent evolution of the solution state, while a contextual bandit approach could condition heuristic selection on features of the current solution, potentially yielding more targeted adaptive behavior. 


%%%%%%%%%%%%%%%%%%%%%%%%%%%%%%%%%%%%%%%%%%%%%%%%%%%%%%%%%%%%%%%%%%%%%%%%
%%
%% The next two lines define the bibliography style to be used, and
%% the bibliography file.
\bibliographystyle{ACM-Reference-Format}
\bibliography{sample-base}


%%
%% If your work has an appendix, this is the place to put it.
\begin{comment}
\appendix

\section{Research Methods}

\subsection{Part One}

Lorem ipsum dolor sit amet, consectetur adipiscing elit. Morbi
malesuada, quam in pulvinar varius, metus nunc fermentum urna, id
sollicitudin purus odio sit amet enim. Aliquam ullamcorper eu ipsum
vel mollis. Curabitur quis dictum nisl. Phasellus vel semper risus, et
lacinia dolor. Integer ultricies commodo sem nec semper.

\subsection{Part Two}

Etiam commodo feugiat nisl pulvinar pellentesque. Etiam auctor sodales
ligula, non varius nibh pulvinar semper. Suspendisse nec lectus non
ipsum convallis congue hendrerit vitae sapien. Donec at laoreet
eros. Vivamus non purus placerat, scelerisque diam eu, cursus
ante. Etiam aliquam tortor auctor efficitur mattis.

\section{Online Resources}

Nam id fermentum dui. Suspendisse sagittis tortor a nulla mollis, in
pulvinar ex pretium. Sed interdum orci quis metus euismod, et sagittis
enim maximus. Vestibulum gravida massa ut felis suscipit
congue. Quisque mattis elit a risus ultrices commodo venenatis eget
dui. Etiam sagittis eleifend elementum.

Nam interdum magna at lectus dignissim, ac dignissim lorem
rhoncus. Maecenas eu arcu ac neque placerat aliquam. Nunc pulvinar
massa et mattis lacinia.
\end{comment}
\appendix
\section{Reproducibility Checklist for JAIR}

Select the answers that apply to your research -- one per item. 

\subsection*{All articles:}

%\hh{revised for stylistic consistency:}
\begin{enumerate}
    \item All claims investigated in this work are clearly stated. 
    [yes]
    \item Clear explanations are given how the work reported substantiates the claims. 
    [yes]
    \item Limitations or technical assumptions are stated clearly and explicitly. 
    [yes]
    \item Conceptual outlines and/or pseudo-code descriptions of the AI methods introduced in this work are provided, and important implementation details are discussed. 
    [yes]
    \item 
    Motivation is provided for all design choices, including algorithms, implementation choices, parameters, data sets and experimental protocols beyond metrics.
    [yes]
\end{enumerate}

\subsection*{Articles containing theoretical contributions:}
Does this paper make theoretical contributions? 
[yes] 

If yes, please complete the list below.

\begin{enumerate}
    \item All assumptions and restrictions are stated clearly and formally. 
    [yes]
    \item All novel claims are stated formally (e.g., in theorem statements). 
    [yes]
    \item Proofs of all non-trivial claims are provided in sufficient detail to permit verification by readers with a reasonable degree of expertise (e.g., that expected from a PhD candidate in the same area of AI). [yes]
    \item
    Complex formalism, such as definitions or proofs, is motivated and explained clearly.
%hh: was:
%Proof sketches or intuitions are given for complex and/or novel results.
    [yes]
    \item 
    The use of mathematical notation and formalism serves the purpose of enhancing clarity and precision; gratuitous use of mathematical formalism (i.e., use that does not enhance clarity or precision) is avoided.
    [yes]
    \item 
    Appropriate citations are given for all non-trivial theoretical tools and techniques. 
    [yes]
\end{enumerate}

\subsection*{Articles reporting on computational experiments:}
Does this paper include computational experiments? [yes]

If yes, please complete the list below.
\begin{enumerate}
    \item 
    All source code required for conducting experiments is included in an online appendix 
    or will be made publicly available upon publication of the paper.
    The online appendix follows best practices for source code readability and documentation as well as for long-term accessibility.
    [yes]
    \item The source code comes with a license that
    allows free usage for reproducibility purposes.
    [yes]
    \item The source code comes with a license that
    allows free usage for research purposes in general.
    [yes]
    \item 
    Raw, unaggregated data from all experiments is included in an online appendix 
    or will be made publicly available upon publication of the paper.
    The online appendix follows best practices for long-term accessibility.
    [yes]
    \item The unaggregated data comes with a license that
    allows free usage for reproducibility purposes.
    [yes]
    \item The unaggregated data comes with a license that
    allows free usage for research purposes in general.
    [yes]
    \item If an algorithm depends on randomness, then the method used for generating random numbers and for setting seeds is described in a way sufficient to allow replication of results. 
    [yes]
    \item The execution environment for experiments, the computing infrastructure (hardware and software) used for running them, is described, including GPU/CPU makes and models; amount of memory (cache and RAM); make and version of operating system; names and versions of relevant software libraries and frameworks. 
    [yes]
    \item 
    The evaluation metrics used in experiments are clearly explained and their choice is explicitly motivated. 
    [yes]
    \item 
    The number of algorithm runs used to compute each result is reported. 
    [yes]
    \item 
    Reported results have not been ``cherry-picked'' by silently ignoring unsuccessful or unsatisfactory experiments. 
    [yes]
    \item 
    Analysis of results goes beyond single-dimensional summaries of performance (e.g., average, median) to include measures of variation, confidence, or other distributional information. 
    [yes]
    \item 
    All (hyper-) parameter settings for 
    the algorithms/methods used in experiments have been reported, along with the rationale or method for determining them. 
    [yes]
    \item 
    The number and range of (hyper-) parameter settings explored prior to conducting final experiments have been indicated, along with the effort spent on (hyper-) parameter optimisation. 
    [yes]
    \item 
    Appropriately chosen statistical hypothesis tests are used to establish statistical significance
    in the presence of noise effects.
    [NA]
\end{enumerate}


\subsection*{Articles using data sets:}
Does this work rely on one or more data sets (possibly obtained from a benchmark generator or similar software artifact)? 
[yes]

If yes, please complete the list below.
\begin{enumerate}
    \item 
    All newly introduced data sets 
    are included in an online appendix 
    or will be made publicly available upon publication of the paper.
    The online appendix follows best practices for long-term accessibility with a license
    that allows free usage for research purposes.
    [yes]
    \item The newly introduced data set comes with a license that
    allows free usage for reproducibility purposes.
    [yes]
    \item The newly introduced data set comes with a license that
    allows free usage for research purposes in general.
    [yes]
    \item All data sets drawn from the literature or other public sources (potentially including authors' own previously published work) are accompanied by appropriate citations.
    [yes]
    \item All data sets drawn from the existing literature (potentially including authors’ own previously published work) are publicly available. [yes]
    %\item All data sets that are not publicly available are described in detail.
    %[yes/partially/no/NA]
    \item All new data sets and data sets that are not publicly available are described in detail, including relevant statistics, the data collection process and annotation process if relevant.
    [yes]
    \item 
    All methods used for preprocessing, augmenting, batching or splitting data sets (e.g., in the context of hold-out or cross-validation)
    are described in detail. [yes]
\end{enumerate}


\end{document}
\endinput
%%
%% End of file `sample-acmlarge.tex'.
